\chapter{Differential and Integral Calculus}\label{ch:calculus}

\begin{center}\itshape
$\dfrac\pi2=\dfrac21\cdot\dfrac23\cdot\dfrac43\cdot\dfrac45\cdot\dfrac65\cdot\dfrac67\cdots$ \quad and \quad $\displaystyle\int_0^\infty e^{-x^2}\dd x=\frac{\sqrt\pi}2$.\\
Both come out of one humble family of integrals, which also completes Stirling's formula.
\end{center}

In 1656 John Wallis wanted a value for $\pi$. In the \emph{Arithmetica Infinitorum} he computed the area under $(1-x^2)^n$ for whole numbers $n$. Each area is a ratio of factorials. At the half-integer $n=\frac12$ the curve is a quarter circle, so the area is $\frac\pi4$. He guessed that the factorial formula still made sense there, and a product for $\frac\pi2$ fell out. He said the step was a guess. Pierre de Fermat, who was sent the book in 1656, answered that a pattern checked on several whole numbers is not a proof. The letters ran through 1657. Wallis printed the exchange in the \emph{Commercium epistolicum} of 1658.

The guess was right. The same areas, written as $W_n=\int_0^{\pi/2}\sin^nt\,\dd t$, satisfy
\[
W_n=\frac{n-1}n\,W_{n-2}\qquad(n\ge2).
\]
That one line produces Wallis's product, the constant $\sqrt{2\pi}$ in Stirling's formula, and the Gaussian integral. Three famous results, one relation between an integral and the previous one.

The other thread is a polynomial with the error kept. Brook Taylor published the expansion in 1715, as Proposition~VII of the \emph{Methodus Incrementorum Directa et Inversa}. He gave no remainder one can bound. The series for sine, cosine and the exponential are instances of his formula; he did not write them out. In 1797, in the \emph{Th\'eorie des fonctions analytiques}, Lagrange took that expansion as the starting point of the differential calculus, and in the same book he gave the remainder a form one can bound. The curve $e^{-x^2}$ enters earlier, in Abraham de Moivre's privately printed pamphlet of 12~November 1733, as an approximation to binomial coefficients. Pierre-Simon Laplace evaluated its integral in 1774.

Every later estimate in this chapter is an integral of a derivative, so the derivative rules and the fundamental theorem come first. Taylor's formula is that comparison with the error written down. L'H\^opital's rule is the same comparison, applied to a quotient. Wallis's integrals are the worked example. Euler's Gamma function is the integral that makes $n!$ a function of a real variable: an infinite product in his letter to Christian Goldbach of 13~October 1729, and an integral in the letter of 8~January 1730. Two short integrations by parts then prove that $e$ and $\pi$ are irrational. Joseph Fourier's argument for $e$ was written down by Jean-Baptiste de Stainville in 1815. Ivan Niven's argument for $\pi$ is one page, from 1947.

\portraits{taylor}{Brook Taylor (1685--1731). Portrait, artist unknown. Public domain, via Wikimedia Commons.}{wallis}{John Wallis (1616--1703). After Godfrey Kneller, 1701. Public domain, via Wikimedia Commons.}

\section{Rules for derivatives}

The derivative and the mean value theorem are already in Chapter~\ref{ch:sequences}. What a working calculation needs next is the rules that build a new derivative out of old ones. We write $f^{(k)}$ for the $k$th derivative, and $f^{(0)}=f$. A function is of class $C^k$ on an open interval when those derivatives through order $k$ exist and are continuous. On a closed interval the derivatives are continuous up to the ends, and an endpoint derivative is one-sided. We ask for that continuity so that, on the closed interval between two points, a continuous derivative attains a finite maximum. Every remainder bound below is a multiple of such a maximum.

\begin{propositionx}[rules of differentiation]
Let $f$ and $g$ be differentiable at an interior point $a$ of their intervals of definition.
\begin{enumerate}[label=(\alph*)]
\item $(f+g)'(a)=f'(a)+g'(a)$ and $(fg)'(a)=f'(a)g(a)+f(a)g'(a)$; if $g(a)\ne0$, $\left(\frac fg\right)'(a)=\frac{f'(a)g(a)-f(a)g'(a)}{g(a)^2}$.
\item \textbf{Chain rule}: if $h$ is differentiable at $f(a)$, then $(h\circ f)'(a)=h'(f(a))\,f'(a)$.
\item \textbf{Inverse functions}: if $f$ is continuous and strictly monotone on an interval $I$, with $a$ interior to $I$ and $f'(a)\ne0$, then $f^{-1}$ is differentiable at $b=f(a)$ with $(f^{-1})'(b)=\frac1{f'(a)}$.
\end{enumerate}
\end{propositionx}
\begin{proof}
(a) and (b) are Proposition~\ref{prop:early-differentiation}, proved in Chapter~\ref{ch:sequences} before the first series that needed them. The product rule is the one to keep in mind: the increment of $fg$ splits into ``$f$ moves, $g$ waits'' and ``$g$ moves, $f$ waits''. Integration by parts is that sentence read backwards.

(c) Strict monotonicity and the intermediate value theorem (Theorem~\ref{thm:ivt}) make $f$ a bijection from $I$ onto the interval $J=f(I)$. The inverse $g=f^{-1}:J\to I$ is monotone. It is continuous. If, say, $g$ increased and jumped at $b$, the one-sided limits $g(b^-)\le g(b)\le g(b^+)$ would exist, because a monotone function on an interval has one-sided limits, and they would not all be equal. Every value strictly between $g(b^-)$ and $g(b^+)$, other than $g(b)$ itself, would be missed. The image $I$ would not be an interval.

Now let $y\in J$, $y\ne b$, and set $x=g(y)$, so $x\ne a$ and $y=f(x)$. Then
\[
\frac{g(y)-g(b)}{y-b}=\frac{x-a}{f(x)-f(a)}.
\]
As $y\to b$, continuity of $g$ forces $x\to a$. The right-hand side tends to $\frac1{f'(a)}$.
\end{proof}

Check the logarithm and the arctangent, both inverses. From $\exp'=\exp$ and $\exp(0)=1$, part~(c) gives $\ln'(x)=\frac1x$ for $x>0$. From $\tan' = 1+\tan^2$ and $\tan'(0)=1$, the same part gives
\[
\arctan'(x)=\frac1{1+x^2},\qquad x\in\R.
\]
The range of $\arctan$ is $\bigl]-\frac\pi2,\frac\pi2\bigr[$, because $\frac\pi2$ is the first positive zero of cosine, as fixed in the reminder below.

\begin{reminder}[what is admitted about the elementary functions]
The volume takes for granted $\exp$, $\ln$, $\cos$, $\sin$ and the number $\pi$, with these properties: $\exp'=\exp$, $\exp(0)=1$, $\exp(a+b)=\exp a\cdot\exp b$, and $\exp$ is a bijection from $\R$ onto $]0,\infty[$ with inverse $\ln$. Also $\sin'=\cos$, $\cos'=-\sin$, $\cos^2+\sin^2=1$, the addition formulas, and $\pi$ is twice the first positive zero of $\cos$. The series $\sum\frac{x^n}{n!}$ is one way to build the exponential. Absolute convergence follows from the ratio test. Exercise~\ref{ex:taylorln}(c) identifies the series with the complex exponential of Chapter~\ref{ch:complex}. Until that exercise, the properties in this reminder are named assumptions, not theorems of this chapter.

From them, $x^\alpha=e^{\alpha\ln x}$ for $x>0$ has derivative $\alpha x^{\alpha-1}$. And $\tan=\frac{\sin}{\cos}$ has derivative $1+\tan^2$, which is the input to the arctangent check above.
\end{reminder}

\section{Antiderivatives, read backwards}

An \emph{antiderivative} of $f$ on an interval is a function $F$ with $F'=f$. The mean value theorem (Theorem~\ref{thm:mvt}(c)) says that two antiderivatives differ by a constant: their difference has derivative zero, so it does not move. Proposition~\ref{prop:inttools} will show that every continuous $f$ has at least one antiderivative, namely $x\mapsto\int_a^x f$. The three techniques below are differentiation rules run backwards. Each proposed answer can be checked by differentiating it.

\begin{examplex}[substitution, parts, partial fractions]
\emph{Substitution} is the chain rule backwards. For $\int\frac{x}{1+x^2}\dd x$, notice $x\dd x=\frac12\dd(1+x^2)$, so the integral is $\frac12\ln(1+x^2)+C$. For the quarter circle, put $x=\sin\theta$ on $[0,1]$, so $\dd x=\cos\theta\dd\theta$ and $\theta$ runs from $0$ to $\frac\pi2$:
\[
\int_0^1\sqrt{1-x^2}\,\dd x=\int_0^{\pi/2}\cos^2\theta\,\dd\theta=\frac\pi4.
\]
The value $\frac\pi4$ is the area you can see in Figure~\ref{fig:quarter}. The same cosine integral is $W_2$, computed again in Proposition~\ref{prop:wallis}.

\emph{Integration by parts} is the product rule backwards. For $\int x\,e^x\dd x$, differentiate the polynomial and integrate the exponential:
\[
xe^x-\int e^x\dd x=(x-1)e^x+C.
\]
For $\int\ln x\dd x$, write the integrand as $1\cdot\ln x$ and differentiate the logarithm:
\[
x\ln x-\int x\cdot\frac1x\dd x=x\ln x-x+C.
\]
The useful habit: differentiate the factor that gets simpler (a polynomial, a logarithm, an inverse trigonometric function), and integrate the factor that stays the same sort of function (an exponential, a sine, a cosine). Figure~\ref{fig:parts} is the picture when one endpoint is at the origin.

\emph{Partial fractions} split a rational function into pieces whose antiderivatives are logarithms or arctangents. For instance
\[
\frac1{x^2-1}=\frac12\left(\frac1{x-1}-\frac1{x+1}\right),
\]
so $\int\frac{\dd x}{x^2-1}=\frac12\ln\left|\frac{x-1}{x+1}\right|+C$, on any interval that does not contain $\pm1$. If the denominator is an irreducible quadratic, complete the square. When $4q>p^2$,
\[
\int\frac{\dd x}{x^2+px+q}
\]
becomes a constant multiple of $\int\frac{\dd u}{1+u^2}=\arctan u$.
\end{examplex}

\begin{figure}[ht]
\centering
\fig{ch06-quarter}
\caption{Left: $y=\sqrt{1-x^2}$ on $[0,1]$ is a quarter circle. The coral region has area $\frac\pi4$, which is the substitution in the example. The rust square has area $1$, so the quarter disc is visibly smaller. Right: the partial products $\prod_{k=1}^p\frac{(2k)^2}{(2k-1)(2k+1)}$ for $p=1,\ldots,8$. They climb toward the dashed line $\frac\pi2$. The first product is $\frac43$. Eight factors have only reached about $1.525$. Exercise~\ref{ex:wallisprod} proves the limit.}
\label{fig:quarter}
\end{figure}

\begin{figure}[ht]
\centering
\fig{ch06-parts}
\caption{A graph with $f(0)=0$. The coral region under the graph and the ochre region to its left fill the rust rectangle, so their areas add to $b\,f(b)$. That identity is integration by parts: $\int_0^b f(x)\dd x+\int_0^{f(b)}f^{-1}(y)\dd y=b\,f(b)$.}
\label{fig:parts}
\end{figure}

\begin{reminder}[improper integrals]
When $f$ is continuous on $[a,\infty[$, define $\int_a^\infty f=\lim_{A\to\infty}\int_a^A f$ when the limit exists, and say the integral \emph{converges}. Likewise $\int_a^b f=\lim_{\varepsilon\to0^+}\int_{a+\varepsilon}^b f$ when $f$ is continuous only on $]a,b]$.

For $f\ge0$ the partial integrals increase with $A$. They converge if and only if they stay bounded. Comparison works as it does for series: $0\le f\le g$ and $\int g<\infty$ give $\int f<\infty$. Thus $\int_1^\infty x^{-\alpha}\dd x$ converges if and only if $\alpha>1$, and $\int_0^1 x^{-\alpha}\dd x$ converges if and only if $\alpha<1$. Also $\int_0^\infty e^{-x}\dd x=1$. The Gaussian integral $\int_0^\infty e^{-x^2}\dd x$ converges because $e^{-x^2}\le e^{-x}$ for $x\ge1$. Its value is $\frac{\sqrt\pi}2$, proved in Exercise~\ref{ex:gaussian} from the Wallis asymptotics below. The exercise does not need polar coordinates.

If $\int|f|$ converges, so does $\int f$. For large $A<B$, the estimate $\bigl|\int_A^B f\bigr|\le\int_A^B|f|$ is as small as you like, so the integrals up to integer endpoints form a Cauchy sequence. The same estimate shows that every real cutoff tends to the same limit. Near a finite improper endpoint the argument is the same, with a sequence of cutoffs.

A concrete conditional example is $\int_0^\infty\frac{\sin x}x\dd x$. On $[1,A]$, integrate by parts with the antiderivative $-\cos x$ of $\sin x$:
\[
\int_1^A\frac{\sin x}x\dd x=\Bigl[-\frac{\cos x}x\Bigr]_1^A-\int_1^A\frac{\cos x}{x^2}\dd x.
\]
The boundary term converges, and the last integral converges absolutely by comparison with $\frac1{x^2}$. At zero, $\frac{\sin x}x$ extends continuously by the value $1$. Absolute convergence fails: on $[k\pi,(k+1)\pi]$ with $k\ge1$, the integral of $\frac{|\sin x|}x$ is at least $\frac2{(k+1)\pi}$, and that series diverges. So the integral converges, and not absolutely. Its value is $\frac\pi2$. Volume~5 evaluates it; the argument uses differentiation under the integral sign, which this chapter does not develop. The convergence is already proved.
\end{reminder}

\section{The integration toolkit}

Chapter~\ref{ch:sequences} fixed the contract we still use. A continuous function on a closed bounded interval has an integral. The integral is linear and additive over adjacent intervals, it respects inequalities, and $\int_a^b c\,\dd t=c(b-a)$. Reversed limits change the sign. Proposition~\ref{prop:early-ftc} proved the fundamental theorem from that contract, together with continuity and the mean value theorem. Volume~7 builds the integral that the contract names, using Heine's theorem, which is already proved. This chapter uses the contract.

\begin{proposition}[integration toolkit]\label{prop:inttools}
Let $f$ be continuous on an interval $I$ and $a\in I$.
\begin{enumerate}[label=(\alph*)]
\item \textbf{Fundamental theorem}: $F(x)=\int_a^xf$ is differentiable with $F'=f$ at every interior point of $I$. Hence $\int_a^bg'=g(b)-g(a)$ for $g$ of class $C^1$.
\item \textbf{Integration by parts}: for $u,v$ of class $C^1$ on $[a,b]$, $\int_a^bu'v=[uv]_a^b-\int_a^buv'$.
\item \textbf{Change of variables}: for $\varphi$ of class $C^1$ on $[a,b]$ with values in $I$,
\[\int_{\varphi(a)}^{\varphi(b)}f(x)\dd x=\int_a^bf(\varphi(t))\,\varphi'(t)\dd t.\]
\item \textbf{Riemann sums}: for $f$ of class $C^1$ on $[0,1]$,
\[\left|\frac1n\sum_{k=1}^nf\left(\frac kn\right)-\int_0^1f\right|\le\frac{\sup|f'|}{n}.\]
\end{enumerate}
\end{proposition}
\begin{proof}
(a) is Proposition~\ref{prop:early-ftc}.

(b) The product $uv$ is an antiderivative of $u'v+uv'$, and that derivative is continuous. Part~(a) gives $\int_a^b(u'v+uv')=[uv]_a^b$. Rearrange.

(c) Fix the lower limit and let the upper limit $b$ vary. Both sides vanish at $b=a$. The derivative of the left side is $f(\varphi(b))\varphi'(b)$, by~(a) and the chain rule. The derivative of the right side is the same, by~(a) applied to the continuous integrand $t\mapsto f(\varphi(t))\varphi'(t)$. Two functions with the same derivative, and the same value at $a$, agree.

(d) On the slab $\bigl[\frac{k-1}n,\frac kn\bigr]$, the mean value theorem gives $\bigl|f(t)-f\bigl(\frac kn\bigr)\bigr|\le\frac{\sup|f'|}n$. Integrate that inequality over a slab of length $\frac1n$:
\[
\left|\int_{(k-1)/n}^{k/n}f-\frac1n f\Bigl(\frac kn\Bigr)\right|\le\frac{\sup|f'|}{n^2}.
\]
Add the $n$ slabs. The errors add up to at most $\frac{\sup|f'|}n$. Figure~\ref{fig:riemannsum}, left, is the picture for $n=8$.
\end{proof}

Part~(d) still holds for a function that is only continuous, except that the error is no longer pinned to $\frac1n$. Heine's theorem (Theorem~\ref{thm:heine}) says: given $\varepsilon>0$, there is a $\delta>0$ such that $|f(t)-f(s)|\le\varepsilon$ whenever $|t-s|\le\delta$. For $n>\frac1\delta$ the same slab computation gives an error of at most $\varepsilon$. The sequences chapter proved Heine's theorem for this conclusion: a continuous function on a closed interval is a limit of finite sums.

\begin{figure}[ht]
\centering
\fig{ch06-riemann}
\caption{Left: right-endpoint Riemann sum for $e^{-x^2}$ on $[0,1]$, with $n=8$ coral bars. Part~(d) bounds the error by $\frac{\sup|f'|}8\approx0.11$. The integral is about $0.747$ and the sum is about $0.706$, so the actual error is about $0.040$. Right: $\sin^nt$ on $\bigl[0,\frac\pi2\bigr]$ for $n=1,2,5,20,100$. A high power is tiny except where the base is near $1$, so the graph collapses toward $t=\frac\pi2$. The width of the bump is of order $\frac1{\sqrt n}$, which is why $W_n\sim\sqrt{\frac\pi{2n}}$ (Proposition~\ref{prop:wallis}).}
\label{fig:riemannsum}
\end{figure}

\section{Taylor's formula, with the error kept}

\begin{think}
Without a calculator, estimate $\sin(0.1)$. The tangent line at $0$ says $0.1$. How far off is that, and which term would you add next? How many terms of the shape $\pm\frac{x^k}{k!}$ do you need before the guaranteed error is smaller than $10^{-6}$?
\end{think}

The fundamental theorem at order zero is already a Taylor formula. If $f$ is $C^1$,
\[
f(x)=f(a)+\int_a^x f'(t)\dd t.
\]
The number $f(a)$ is the constant polynomial. The integral is the error. Integrate that error by parts, again and again, and each step peels off one more term $\frac{f^{(k)}(a)}{k!}(x-a)^k$. What remains is still an integral. That integral is the remainder.

\begin{theorem}[Taylor formulas]\label{thm:taylor}
Let $f$ be of class $C^{n+1}$ on an interval containing $a$ and $x$, and $T_n(x)=\sum_{k=0}^n\frac{f^{(k)}(a)}{k!}(x-a)^k$.
\begin{enumerate}[label=(\alph*)]
\item \textbf{Integral remainder}:
\[f(x)=T_n(x)+\int_a^x\frac{(x-t)^n}{n!}f^{(n+1)}(t)\dd t.\]
\item \textbf{Taylor--Lagrange inequality}:
\[|f(x)-T_n(x)|\le\frac{|x-a|^{n+1}}{(n+1)!}\sup|f^{(n+1)}|,\]
the supremum over $J=[\min(a,x),\max(a,x)]$.
\item \textbf{Taylor--Young}: if $f$ is only $C^n$, $f(x)=T_n(x)+o\bigl((x-a)^n\bigr)$ as $x\to a$.
\end{enumerate}
\end{theorem}
\begin{proof}
(a) Proceed by induction on $n$. The case $n=0$ is Proposition~\ref{prop:inttools}(a). Suppose the formula holds at order $n-1$, for a function of class $C^{n+1}$. The remainder at that order is
\[
\int_a^x\frac{(x-t)^{n-1}}{(n-1)!}f^{(n)}(t)\dd t.
\]
Integrate by parts. Take
\[
u'=\frac{(x-t)^{n-1}}{(n-1)!},\qquad u=-\frac{(x-t)^n}{n!},\qquad v=f^{(n)},
\]
so $v'=f^{(n+1)}$. The product $uv$ vanishes at $t=x$, because $u$ does. At $t=a$,
\[
u(a)v(a)=-\frac{(x-a)^n}{n!}f^{(n)}(a).
\]
Proposition~\ref{prop:inttools}(b) therefore turns the integral into
\[
\frac{(x-a)^n}{n!}f^{(n)}(a)+\int_a^x\frac{(x-t)^n}{n!}f^{(n+1)}(t)\dd t.
\]
The first summand is the next Taylor term. The integral is the remainder at order $n$.

(b) Let $M=\sup_J|f^{(n+1)}|$. The integrand in~(a) has absolute value at most $M\frac{|x-t|^n}{n!}$. The comparison integral is elementary: the substitution $u=x-t$ gives
\[
\left|\int_a^x\frac{|x-t|^n}{n!}\dd t\right|=\frac{|x-a|^{n+1}}{(n+1)!}.
\]
Multiply by $M$.

(c) For $n=0$ the claim is continuity of $f$. For $n\ge1$, the function is $C^n$, so part~(a) applies at order $n-1$. Subtract the integral that produces the last Taylor term,
\[
\frac{f^{(n)}(a)}{n!}(x-a)^n=\int_a^x\frac{(x-t)^{n-1}}{(n-1)!}f^{(n)}(a)\dd t.
\]
The difference is an integral of $f^{(n)}(t)-f^{(n)}(a)$. The same comparison as in~(b) gives
\[
|f(x)-T_n(x)|\le\frac{|x-a|^n}{n!}\sup_{t\in J}\bigl|f^{(n)}(t)-f^{(n)}(a)\bigr|.
\]
The supremum tends to $0$ as $x\to a$, because $f^{(n)}$ is continuous at $a$.
\end{proof}

Lagrange's 1797 form names a point $c$ between $a$ and $x$ and writes the error as $f^{(n+1)}(c)\frac{(x-a)^{n+1}}{(n+1)!}$. The inequality in~(b) is the consequence used below. It needs a bound on the derivative, which is what a numerical estimate has, and the proof does not have to locate $c$.

\begin{examplex}[$\sin(0.1)$, as promised]
Every derivative of sine is $\pm\sin$ or $\pm\cos$, hence bounded by $1$. Theorem~\ref{thm:taylor}(b) at $a=0$ gives $|\sin x-T_n(x)|\le\frac{|x|^{n+1}}{(n+1)!}$. The tangent line $T_1(x)=x$ approximates $\sin(0.1)$ by $0.1$, with error at most $\frac{0.01}2=0.005$. The actual gap is about $0.000167$. The bound is crude because the second derivative of sine vanishes at $0$, so $T_1=T_2$. Using $n=2$ replaces the bound by $\frac{0.001}6\approx1.67\times10^{-4}$.

The next polynomial is $T_3(x)=x-\frac{x^3}6$. At $x=0.1$,
\[
T_3(0.1)=0.1-\frac{0.001}6=0.0998333\ldots.
\]
Again $T_3=T_4$, because the fourth derivative vanishes at $0$. The bound at order $4$ is $\frac{10^{-5}}{120}<10^{-7}$. So two nonzero terms already guarantee six correct decimals. The value begins $0.0998334\ldots$. Figure~\ref{fig:zoom} is this window; Figure~\ref{fig:taylor}, left, is the same story on a much larger one.
\end{examplex}

The series of sine converges to $\sin x$ for every real $x$, because $\frac{|x|^{n+1}}{(n+1)!}\to0$. At $x=\pi$ the bound stays large until $n$ is well past $\pi$. A low-degree polynomial is most useful near its centre. Figure~\ref{fig:taylorexp} shows two other fates. The polynomials of the exponential converge at every fixed $x$. The polynomials of $\ln(1+x)$ converge on $]-1,1]$ and nowhere else: past $x=1$ the terms do not tend to $0$.

\begin{figure}[ht]
\centering
\fig{ch06-zoom}
\caption{Sine on $[0,1]$, in charcoal, against $T_1$ (rust), $T_3$ (ochre) and $T_5$ (brown). The tangent pulls away as $x$ grows. At the marked point $x=0.1$ the gap is smaller than the line width: $T_3$ is already within $10^{-7}$, by the Lagrange bound.}
\label{fig:zoom}
\end{figure}

\begin{figure}[ht]
\centering
\fig{ch06-taylor}
\caption{Left: $\sin$ and its Taylor polynomials $T_1,T_5,T_9,T_{13}$ at $0$. At $x=\pi$ they give $3.14$, $0.524$, $0.0069$ and $0.00002$, closing in on $\sin\pi=0$. Right: the sandwich of Exercise~\ref{ex:gaussian} for $n=3$. The three areas are $\sqrt3\,W_7\approx0.792$, $\frac{\sqrt\pi}2\approx0.886$ and $\sqrt3\,W_4\approx1.020$. As $n$ grows, both bounding areas close in on the Gaussian integral.}
\label{fig:taylor}
\end{figure}

\begin{figure}[ht]
\centering
\fig{ch06-taylorexp}
\caption{Left: $e^x$ and $T_1,\dots,T_4$ at $0$. For each fixed $x$ the polynomials converge to $e^x$. When $x<0$, the next degree need not be the closer one. Right: $\ln(1+x)$ and $T_2,T_5,T_{10},T_{20}$ at $0$. They converge on $]-1,1[$, and still at $x=1$, slowly (Exercise~\ref{ex:taylorln}). For a fixed $x>1$ the terms fail to tend to $0$, so the series diverges. The radius of that series is $1$. Volume~5 studies radii in general.}
\label{fig:taylorexp}
\end{figure}

\section{L'H\^opital's rule}

The first printed calculus textbook, the \emph{Analyse des infiniment petits} of 1696, carries the name of the Marquis de l'H\^opital, and so does its best-known rule. On 17~March 1694 he wrote to Johann Bernoulli offering a retainer of 300~livres a year, from the first of January, in exchange for Bernoulli's time, his discoveries, and his silence. Bernoulli accepted. He objected after the marquis died. In 1921 a manuscript of the course Bernoulli had given turned up; the book follows it closely. The rule for the form $\frac00$ is part of what that course contained. The proof below does not depend on the contract. It is Cauchy's mean value theorem, which is Rolle's theorem applied to a mixture of $f$ and $g$.

\begin{theorem}[l'H\^opital's rule]\label{thm:lhopital}
\begin{enumerate}[label=(\alph*)]
\item Let $a<b$ and let $f,g$ be continuous on $[a,b]$ and differentiable on $]a,b[$. Then, for some $c\in\,]a,b[$,
\[
\bigl(f(b)-f(a)\bigr)g'(c)=\bigl(g(b)-g(a)\bigr)f'(c).
\]
\item Let $f,g$ be differentiable on $]a,a+h[$ with $g'\ne0$ there, and suppose $f(x)\to0$ and $g(x)\to0$ as $x\to a^+$. If $\frac{f'(x)}{g'(x)}\to\ell$ as $x\to a^+$ (with $\ell\in\R$ or $\ell=\pm\infty$), then $\frac{f(x)}{g(x)}\to\ell$ as well. The same holds for $x\to a^-$, for two-sided limits, and for $x\to\pm\infty$.
\end{enumerate}
\end{theorem}
\begin{proof}
(a) Define $\varphi(x)=\bigl(f(b)-f(a)\bigr)g(x)-\bigl(g(b)-g(a)\bigr)f(x)$. Then $\varphi$ is continuous on $[a,b]$ and differentiable inside. At both endpoints,
\[
\varphi(a)=\varphi(b)=f(b)g(a)-g(b)f(a).
\]
Rolle's theorem (Theorem~\ref{thm:mvt}) gives a point $c$ with $\varphi'(c)=0$, which is the identity claimed.

(b) Extend $f$ and $g$ by $f(a)=g(a)=0$. They are then continuous on $[a,a+h[$. Also $g(x)\ne0$ for $x\in\,]a,a+h[$: a zero would give, by Rolle, a zero of $g'$. For such an $x$, part~(a) on $[a,x]$ produces $c_x\in\,]a,x[$ with
\[
\frac{f(x)}{g(x)}=\frac{f'(c_x)}{g'(c_x)}.
\]
As $x\to a^+$, the point $c_x$ is squeezed to $a^+$ as well, so the right-hand side tends to $\ell$. Left-hand and two-sided limits are the same argument on the corresponding side.

For $x\to+\infty$, set $x=\frac1t$ and $F(t)=f\bigl(\frac1t\bigr)$, $G(t)=g\bigl(\frac1t\bigr)$. As $t\to0^+$, both $F$ and $G$ tend to $0$. The chain rule gives $F'(t)=-\frac1{t^2}f'\bigl(\frac1t\bigr)$ and the same factor in $G'(t)$. The factors cancel in the quotient, for $t>0$, and $\frac{F'}{G'}$ tends to $\ell$. The case already proved yields $\frac{F}{G}\to\ell$. The case $x\to-\infty$ is symmetric.
\end{proof}

\begin{examplex}[four quotients, and one warning]
$\frac{\sin x}x\to1$ as $x\to0$, because the derivatives give $\frac{\cos x}1\to1$. For $\frac{1-\cos x}{x^2}$, one differentiation produces $\frac{\sin x}{2x}$, and a second produces $\frac{\cos x}2\to\frac12$. The rule may be repeated as long as the hypotheses hold at each stage.

$\frac{\ln x}x\to0$ as $x\to\infty$. The quotient of derivatives is $\frac{1/x}{1}=\frac1x\to0$. This is the form $\frac\infty\infty$. Exercise~\ref{ex:lhopital} gives it a proof of the same shape. There is also an elementary bound. Chapter~\ref{ch:sequences} has $\ln n\le4n^{1/4}$. For real $x\ge1$, $\ln x\le\ln(\lfloor x\rfloor+1)\le4(\lfloor x\rfloor+1)^{1/4}$, so $\frac{\ln x}x\to0$ as well.

Taylor--Young is often shorter. From $e^x=1+x+\frac{x^2}2+o(x^2)$,
\[
\frac{e^x-1-x}{x^2}=\frac{\frac{x^2}2+o(x^2)}{x^2}\to\frac12
\]
in one line. L'H\^opital needs two differentiations for the same fact.

The warning: if $\frac{f'}{g'}$ has no limit, the rule is silent. The quotient $\frac fg$ may still converge. As $x\to\infty$, $\frac{x+\sin x}x\to1$, while $\frac{1+\cos x}1$ oscillates and has no limit. A rule that does not apply is not a counterexample to the rule.
\end{examplex}

\section{Wallis integrals}

For $n\in\N$ let $W_n=\displaystyle\int_0^{\pi/2}\sin^nt\dd t$. The substitution $t\mapsto\frac\pi2-t$ gives the same value with cosine in place of sine. A high power of a number in $[0,1)$ is small. On $\bigl[0,\frac\pi2\bigr)$, sine is such a number, except near $\frac\pi2$, where it equals $1$. The integral $W_n$ is the area of that last bump. Figure~\ref{fig:riemannsum}, right, shows the bump getting thinner.

\begin{think}
Sketch $\sin^nt$ on $\bigl[0,\frac\pi2\bigr]$ for $n=1$, $n=10$ and $n=100$. The graphs collapse toward $\frac\pi2$, so $W_n\to0$. How fast? Near the right endpoint, $\sin\bigl(\frac\pi2-s\bigr)=\cos s$, and $\cos s$ is close to $e^{-s^2/2}$ for small $s$. The region where the integrand is not small should have width of order $\frac1{\sqrt n}$.
\end{think}

\begin{proposition}\label{prop:wallis}
\begin{enumerate}[label=(\roman*)]
\item $W_0=\frac\pi2$, $W_1=1$, and $W_n=\frac{n-1}nW_{n-2}$ for $n\ge2$.
\item $W_{2p}=\dfrac{(2p)!}{4^p(p!)^2}\cdot\dfrac\pi2$ and $W_{2p+1}=\dfrac{4^p(p!)^2}{(2p+1)!}$.
\item $(W_n)$ is decreasing and positive, $nW_nW_{n-1}=\frac\pi2$, and $W_n\sim\sqrt{\dfrac\pi{2n}}$.
\end{enumerate}
\end{proposition}
\begin{proof}
(i) The two starting values are immediate: $W_0$ is the length of the interval, and $W_1=\int_0^{\pi/2}\sin t\dd t=1$. For $n\ge2$, integrate by parts with $u'=\sin t$ and $v=\sin^{n-1}t$, so $u=-\cos t$ and $v'=(n-1)\sin^{n-2}t\cos t$. The boundary term is
\[
\bigl[-\cos t\sin^{n-1}t\bigr]_0^{\pi/2}=0,
\]
because $\cos\frac\pi2=0$ and $\sin0=0$. What remains is
\[
W_n=(n-1)\int_0^{\pi/2}\cos^2t\,\sin^{n-2}t\dd t.
\]
Write $\cos^2=1-\sin^2$. The integral splits into $W_{n-2}-W_n$, so $W_n=(n-1)(W_{n-2}-W_n)$. Add $(n-1)W_n$ to both sides and divide by $n$.

Check the first two steps by hand. $W_2=\frac12 W_0=\frac\pi4$, which is the quarter-circle integral of the example above. $W_3=\frac23 W_1=\frac23$.

(ii) Iterate the recurrence. For the even case,
\[
W_{2p}=\frac{(2p-1)(2p-3)\cdots1}{(2p)(2p-2)\cdots2}\,W_0.
\]
The denominator is $2^pp!$. Multiply numerator and denominator by that same product. The numerator becomes $(2p)!$, and the denominator becomes $4^p(p!)^2$. The odd case is the same iteration, starting from $W_1=1$.

(iii) On $\bigl[0,\frac\pi2\bigr]$ one has $0\le\sin^nt\le\sin^{n-1}t$, so $(W_n)$ decreases. The integrand is positive on a subinterval, so $W_n>0$. From~(i),
\[
nW_nW_{n-1}=(n-1)W_{n-1}W_{n-2}.
\]
The product $nW_nW_{n-1}$ does not depend on $n$. Its value at $n=1$ is $W_1W_0=\frac\pi2$.

Finally $\frac{W_n}{W_{n-2}}=\frac{n-1}n$, and the sequence decreases, so
\[
\frac{n-1}n=\frac{W_n}{W_{n-2}}\le\frac{W_n}{W_{n-1}}\le1.
\]
Thus $W_n\sim W_{n-1}$. Square the first of them: $W_n^2\sim nW_nW_{n-1}\cdot\frac1n=\frac\pi{2n}$.
\end{proof}

\begin{think}
Compute $\dfrac{10!}{\sqrt{10}\,(10/e)^{10}}$. You should get about $2.5276$. Which constant is close, and which correction of size $\frac1{12n}$ brings it closer still?
\end{think}

\begin{theorem}[Stirling's constant]\label{thm:stirlingC}
The constant of Exercise~\ref{ex:stirling} is $C=\sqrt{2\pi}$, so $n!\sim\sqrt{2\pi n}\left(\frac ne\right)^n$.
\end{theorem}
\begin{proof}
Exercise~\ref{ex:stirling}, in Chapter~\ref{ch:sequences}, gives $m!\sim C\sqrt m\bigl(\frac me\bigr)^m$ for some constant $C>0$. Insert that asymptotic for $(2p)!$ and for $(p!)^2$ into part~(ii):
\begin{align*}
W_{2p}
&=\frac\pi2\cdot\frac{(2p)!}{4^p(p!)^2}\\
&\sim\frac\pi2\cdot\frac{C\sqrt{2p}\,(2p/e)^{2p}}{4^p\,C^2 p\,(p/e)^{2p}}
=\frac\pi2\cdot\frac{\sqrt{2p}}{C\,p}\cdot\frac{(2p)^{2p}}{4^p\,p^{2p}}.
\end{align*}
The powers cancel. Indeed $(2p)^{2p}=2^{2p}p^{2p}$ and $4^p=(2^2)^p=2^{2p}$, so the last fraction equals $1$. Also $\frac{\sqrt{2p}}{p}=\frac{\sqrt2}{\sqrt p}$. Therefore
\[
W_{2p}\sim\frac\pi2\cdot\frac{\sqrt2}{C\sqrt p}.
\]
Part~(iii) gives the other expression $W_{2p}\sim\sqrt{\frac\pi{4p}}=\frac{\sqrt\pi}{2\sqrt p}$. Match the coefficients of $\frac1{\sqrt p}$:
\[
\frac{\pi\sqrt2}{2C}=\frac{\sqrt\pi}2.
\]
Hence $C=\sqrt{2\pi}$.
\end{proof}

In 1730 de Moivre had the shape $n!\sim B\sqrt n\,(n/e)^n$ and computed $B$ from a slowly convergent series. The partial sum he kept gave $B\approx2.5074$. Stirling identified $B=\sqrt{2\pi}\approx2.5066$, and de Moivre printed the identification with credit to Stirling. The box above, at $n=10$, gives $2.5276$. The one-term correction $\sqrt{2\pi}\bigl(1+\frac1{120}\bigr)$ equals $2.5275$, within $10^{-4}$ of that ratio. Proposition~\ref{prop:stirlingnext} shows where the $\frac1{12n}$ comes from.

\portraits{demoivre}{Abraham de Moivre (1667--1754). Portrait, 1736, artist unknown. Public domain, via Wikimedia Commons.}{euler}{Leonhard Euler (1707--1783). Portrait by Jakob Emanuel Handmann, 1753, Kunstmuseum Basel. Public domain, via Wikimedia Commons.}

\section{The Gamma function}

Half a factorial sounds like a joke. The integral below takes it seriously. On 13~October 1729 Euler wrote to Goldbach with an infinite product that equals $n!$ at the positive integers and makes sense for other positive reals. He had already seen that the product at $\frac12$ is $\frac{\sqrt\pi}2$, by comparison with Wallis. On 8~January 1730 he sent an integral. In modern notation it is $\int_0^1(-\ln x)^n\dd x$. The substitution $t=-\ln x$ turns it into $\int_0^\infty t^n e^{-t}\dd t$. The function below is that integral with the exponent shifted by one, so that the value at a positive integer is a factorial.

\begin{proposition}[the Gamma function]\label{prop:gammafn}
For $x>0$ the integral $\Gamma(x)=\displaystyle\int_0^\infty t^{x-1}e^{-t}\dd t$ converges, and
\begin{enumerate}[label=(\alph*)]
\item $\Gamma(x+1)=x\,\Gamma(x)$ for $x>0$; hence $\Gamma(n+1)=n!$ for $n\in\N$.
\item $\Gamma\bigl(\frac12\bigr)=\sqrt\pi$, and more generally $\Gamma\bigl(n+\frac12\bigr)=\dfrac{(2n)!}{4^n n!}\sqrt\pi$.
\end{enumerate}
\end{proposition}
\begin{proof}
Near $0$, one has $t^{x-1}e^{-t}\le t^{x-1}$, and $\int_0^1 t^{x-1}\dd t=\frac1x$ converges because $x-1>-1$. Near infinity, the bound $\ln n\le4n^{1/4}$ from Chapter~\ref{ch:sequences} gives $\ln t\le4t^{1/4}$ for real $t\ge1$, by the same estimate $\ln u\le u-1$. Hence $t^{x-1}e^{-t/2}\to0$, so $t^{x-1}e^{-t}\le e^{-t/2}$ for all large $t$, and $\int_1^\infty e^{-t/2}\dd t$ converges.

(a) On $[\varepsilon,A]$, integrate by parts with $u=t^x$ and $v'=e^{-t}$, so $v=-e^{-t}$:
\[
\int_\varepsilon^A t^x e^{-t}\dd t=\bigl[-t^x e^{-t}\bigr]_\varepsilon^A+x\int_\varepsilon^A t^{x-1}e^{-t}\dd t.
\]
As $\varepsilon\to0^+$ and $A\to\infty$, the boundary term tends to $0$ for $x>0$: at infinity the exponential beats the power, and at zero $t^x\to0$. The two integrals tend to $\Gamma(x+1)$ and $\Gamma(x)$. Then $\Gamma(1)=\int_0^\infty e^{-t}\dd t=1$, and induction gives $\Gamma(n+1)=n!$.

(b) Substitute $t=u^2$, so $\dd t=2u\dd u$:
\[
\Gamma\Bigl(\frac12\Bigr)=\int_0^\infty t^{-1/2}e^{-t}\dd t=2\int_0^\infty e^{-u^2}\dd u.
\]
Exercise~\ref{ex:gaussian} says the last integral equals $\frac{\sqrt\pi}2$, once Proposition~\ref{prop:wallis} is known. Hence $\Gamma\bigl(\frac12\bigr)=\sqrt\pi$. Apply~(a) repeatedly:
\[
\Gamma\Bigl(n+\frac12\Bigr)=\Bigl(n-\frac12\Bigr)\Bigl(n-\frac32\Bigr)\cdots\frac12\cdot\sqrt\pi=\frac{(2n-1)(2n-3)\cdots1}{2^n}\sqrt\pi.
\]
The double product of odd numbers is $\frac{(2n)!}{2^n n!}$, and the powers of $2$ become $4^n$.
\end{proof}

So $\bigl(\frac12\bigr)!$, if the symbol means $\Gamma\bigl(\frac32\bigr)$, equals $\frac{\sqrt\pi}2\approx0.886$. That is the value Wallis had interpolated in 1656. Figure~\ref{fig:gammafn} shows the graph: the factorials sit on it, there is a minimum near $x=1.46$, and the integral blows up as $x\to0^+$.

Many functions satisfy $f(1)=1$ and $f(x+1)=xf(x)$. Chapter~\ref{ch:inequalities}, Proposition~\ref{prop:gamma-convexity}, proves that $\ln\Gamma$ is convex. Volume~5, Chapter~5, proves the theorem of Bohr and Mollerup from 1922: on $]0,\infty[$, the function $\Gamma$ is the only one with $\Gamma(1)=1$, the functional equation $\Gamma(x+1)=x\Gamma(x)$, and $\ln\Gamma$ convex. The same volume, Chapter~9, extends $\Gamma$ to the complex plane, and Chapter~10 puts it next to $\zeta(s)$ in Riemann's functional equation. Uniqueness is not proved in this chapter.

\begin{figure}[ht]
\centering
\fig{ch06-gammafn}
\caption{Left: $\Gamma$ on $]0,4.5]$, with the points $(n+1,n!)$ marked. The minimum on the positive reals is near $x=1.4616$, where $\Gamma\approx0.8856$. At $x=\frac12$ the value is $\sqrt\pi\approx1.772$. Right: the integrands $t^{x-1}e^{-t}$ for $x=\frac12,1,2,3,4$. The area under each curve is $\Gamma(x)$. As $x$ grows, the bump moves right and its area grows like $(x-1)!$.}
\label{fig:gammafn}
\end{figure}

\section{Numerical integration, and one further term of Stirling}

Most integrals have no elementary antiderivative. A Riemann sum still computes them. Proposition~\ref{prop:inttools}(d) says that $n$ right-hand rectangles leave an error of order $\frac1n$. Two small changes improve the power of $n$: sample at the midpoint of each slab, or replace the rectangle by a trapezoid.

\begin{proposition}[midpoint and trapezoid rules]\label{prop:quadrature}
Let $f$ be of class $C^2$ on $[a,b]$ with $|f''|\le M$, let $h=\frac{b-a}n$ and $x_k=a+kh$. Then
\[
\begin{gathered}
\left|\int_a^bf-h\sum_{k=1}^nf\!\left(x_k-\tfrac h2\right)\right|\le\frac{M(b-a)^3}{24n^2},\\
\left|\int_a^bf-h\sum_{k=1}^n\frac{f(x_{k-1})+f(x_k)}2\right|\le\frac{M(b-a)^3}{12n^2}.
\end{gathered}
\]
\end{proposition}
\begin{proof}
It is enough to bound the error on one slab of length $h$, by $\frac{Mh^3}{24}$ for the midpoint and by $\frac{Mh^3}{12}$ for the trapezoid, and then add the $n$ contributions.

Midpoint. Let the slab be $\bigl[c-\frac h2,c+\frac h2\bigr]$. Taylor's inequality at $c$, to order $1$, gives $f(x)=f(c)+f'(c)(x-c)+r(x)$ with $|r(x)|\le\frac M2(x-c)^2$. Integrate. The linear term is an odd function on a symmetric interval, so its integral is $0$. The comparison integral is
\[
\int_{-h/2}^{h/2}u^2\dd u=2\int_0^{h/2}u^2\dd u=2\cdot\frac1{3}\Bigl(\frac h2\Bigr)^3=\frac{h^3}{12}.
\]
Therefore $\bigl|\int f-hf(c)\bigr|=\bigl|\int r\bigr|\le\frac M2\cdot\frac{h^3}{12}=\frac{Mh^3}{24}$.

Trapezoid. Let the slab be $[\alpha,\beta]$ with $\beta-\alpha=h$, and set
\[
E=\int_\alpha^\beta f-h\frac{f(\alpha)+f(\beta)}2.
\]
The quadratic $p(x)=\frac{(x-\alpha)(\beta-x)}2$ vanishes at both ends, and $p''=-1$. Also $p'(\alpha)=\frac h2$, $p'(\beta)=-\frac h2$, and $p\ge0$ on the slab. Integrate $f=-f\,p''$ by parts twice:
\begin{align*}
\int_\alpha^\beta f
&=-\bigl[fp'\bigr]_\alpha^\beta+\int_\alpha^\beta f'p'\\
&=h\frac{f(\alpha)+f(\beta)}2+\bigl[f'p\bigr]_\alpha^\beta-\int_\alpha^\beta f''p.
\end{align*}
The term $\bigl[f'p\bigr]_\alpha^\beta$ vanishes. Hence $E=-\int_\alpha^\beta f''p$. For the integral of $p$, set $t=x-\alpha$. Then $p=\frac{t(h-t)}2$ and
\[
\int_0^h\frac{t(h-t)}2\dd t=\frac12\Bigl[\frac{ht^2}2-\frac{t^3}3\Bigr]_0^h=\frac12\Bigl(\frac{h^3}2-\frac{h^3}3\Bigr)=\frac{h^3}{12}.
\]
Since $p\ge0$, one has $|E|\le M\frac{h^3}{12}$.
\end{proof}

For a convex function the chord lies above the graph, so the trapezoid overestimates. The midpoint rectangle has the same area as the trapezoid under the tangent at the midpoint, and that tangent lies below the graph, so the midpoint underestimates. The two errors point in opposite directions. The combination $\frac{2M+T}3$ is Simpson's rule. If $f$ is of class $C^4$, its error is of order $\frac1{n^4}$. The proof is the same integration by parts, carried two orders further; it is not written here. Figure~\ref{fig:quadrature} records the four slopes for $\int_0^1 e^{-x^2}\dd x$, an integrand that does have four continuous derivatives.

\begin{figure}[ht]
\centering
\fig{ch06-quadrature}
\caption{Left: on one slab, for a convex graph, the trapezoid overestimates and the coral midpoint rectangle underestimates. The midpoint rectangle has the same area as the trapezoid under the tangent. Right: errors of the rectangle, trapezoid, midpoint and Simpson rules for $\int_0^1 e^{-x^2}\dd x$, against $n$, on logarithmic scales. The slopes are $-1$, $-2$, $-2$ and $-4$. The proposition supplies the midpoint and trapezoid bounds. The Simpson slope uses four continuous derivatives.}
\label{fig:quadrature}
\end{figure}

\begin{proposition}[Stirling, one term further]\label{prop:stirlingnext}
\[n!=\sqrt{2\pi n}\left(\frac ne\right)^n\left(1+\frac1{12n}+O\!\left(\frac1{n^2}\right)\right).\]
\end{proposition}
\begin{proof}
Let $d_n=\ln n!-\bigl(n+\frac12\bigr)\ln n+n$, as in Exercise~\ref{ex:stirling}. Theorem~\ref{thm:stirlingC} says $d_n\to d=\ln\sqrt{2\pi}$. The difference of two consecutive terms collapses:
\[
d_n-d_{n+1}=\Bigl(n+\frac12\Bigr)\ln\Bigl(1+\frac1n\Bigr)-1.
\]
Expand $\ln(1+x)=x-\frac{x^2}2+\frac{x^3}3-\frac{x^4}4+O(x^5)$ at $x=\frac1n$. Split the factor $n+\frac12$ into $n$ and $\frac12$. The first piece produces $1-\frac1{2n}+\frac1{3n^2}-\frac1{4n^3}$. The second produces $\frac1{2n}-\frac1{4n^2}+\frac1{6n^3}$. Add the two pieces:
\begin{align*}
\Bigl(n+\frac12\Bigr)\ln\Bigl(1+\frac1n\Bigr)
&=1-\frac1{2n}+\frac1{3n^2}-\frac1{4n^3}\\
&\qquad+\frac1{2n}-\frac1{4n^2}+\frac1{6n^3}+O\Bigl(\frac1{n^4}\Bigr)\\
&=1+\frac1{12n^2}-\frac1{12n^3}+O\Bigl(\frac1{n^4}\Bigr).
\end{align*}
The coefficient of $\frac1n$ cancels, and $\frac13-\frac14=\frac1{12}$, while $-\frac14+\frac16=-\frac1{12}$. Subtract $1$:
\[
d_n-d_{n+1}=\frac1{12n^2}-\frac1{12n^3}+O\Bigl(\frac1{n^4}\Bigr).
\]
Sum from $n$ onward. The series of differences is a telescoping series, so
\[
d_n-d=\sum_{k\ge n}(d_k-d_{k+1})=\frac1{12}\sum_{k\ge n}\frac1{k^2}+O\Bigl(\sum_{k\ge n}\frac1{k^3}\Bigr).
\]
Compare $\frac1{k^2}$ with telescoping fractions: $\frac1{k(k+1)}\le\frac1{k^2}\le\frac1{(k-1)k}$ for $k\ge2$. Summing from $k=n$ gives
\[
\frac1n\le\sum_{k\ge n}\frac1{k^2}\le\frac1{n-1}=\frac1n+O\Bigl(\frac1{n^2}\Bigr).
\]
For the cubic sum, $\frac1{k^3}\le\frac1n\cdot\frac1{(k-1)k}$ when $k\ge n\ge2$, so $\sum_{k\ge n}\frac1{k^3}\le\frac1{n(n-1)}=O\bigl(\frac1{n^2}\bigr)$. Therefore $d_n-d=\frac1{12n}+O\bigl(\frac1{n^2}\bigr)$. Exponentiate. The expansion $e^u=1+u+O(u^2)$, at a $u$ of size $\frac1n$, gives
\[
\frac{n!}{\sqrt{2\pi n}\,(n/e)^n}=e^{d_n-d}=1+\frac1{12n}+O\Bigl(\frac1{n^2}\Bigr).
\]
\end{proof}

At $n=10$ the factor $1+\frac1{120}$ turns $\sqrt{2\pi}\approx2.5066$ into $2.5275$, and the ratio in the pause above is $2.5276$. The full expansion continues
\[
1+\frac1{12n}+\frac1{288n^2}-\frac{139}{51840n^3}-\cdots.
\]
The series diverges for every fixed $n$: the coefficients eventually grow like factorials. Its first terms are nevertheless extremely accurate. That is an asymptotic series. Volume~5, Chapter~11, proves Stirling's series with a remainder. Volume~6 spends expansions of this kind on the primes. The divergence, and the general term, are not proved here.

\portraits{leibniz}{Gottfried Wilhelm Leibniz (1646--1716). Portrait by Christoph Bernhard Francke, c.\,1700. Public domain, via Wikimedia Commons.}{riemann}{Bernhard Riemann (1826--1866). Photograph, 1863. Public domain, via Wikimedia Commons.}

\section{Workshop: a bound you can compute}

Taylor's formula has two parts. The polynomial is something you can calculate. The remainder is the strap that keeps the calculation attached to the function. A polynomial with the remainder thrown away is a parachute with the straps left at home.

\begin{examplex}[how many digits of $e$?]
For $0\le x\le1$, Theorem~\ref{thm:taylor}(b) gives
\[
\left|e^x-\sum_{k=0}^{n}\frac{x^k}{k!}\right|\le\frac{e}{(n+1)!},
\]
because every derivative of the exponential is the exponential, and $e^x\le e$ on this interval. At $x=1$ the right-hand side tends to $0$, since factorials beat any fixed number. Therefore $e=\sum_{k\ge0}\frac1{k!}$. That identity still has the symbol $e$ in it. A numerical bound needs a number already in hand. For $k\ge2$, one has $k!\ge2^{k-1}$, and the inequality is strict for $k\ge3$, so
\[
e=\sum_{k\ge0}\frac1{k!}<1+1+\sum_{k\ge2}2^{-(k-1)}=3.
\]
At $n=9$ the error is at most $\frac3{10!}<10^{-6}$. This is an absolute error. It does not by itself decide a rounding boundary that sits exactly on a digit $5$. Figure~\ref{fig:review06} draws the same strap at order $2$, with the majorant $3$ in place of $e$.
\end{examplex}

\begin{figure}[ht]\centering\fig{review06-taylor}
\caption{Charcoal curve: $e^x$. Dashed: $T_2(x)=1+x+\frac{x^2}2$, on $[0,1]$. Dotted: the upper curve $T_2(x)+\frac{3x^3}6$. The integral remainder puts the exponential between the two polynomial graphs on this interval.}
\label{fig:review06}
\end{figure}

\subsection*{An integral with no elementary antiderivative}
The job is to approximate $I=\int_0^1 e^{-x^2}\dd x$ to within $10^{-4}$, using rational numbers only. There is no elementary antiderivative to evaluate. The series of the exponential, stopped at a definite order, is enough.

Expand $e^{-u}$ at $u=0$ and then put $u=x^2$. On $[0,1]$ every derivative of $e^{-u}$ has absolute value at most $1$, so Theorem~\ref{thm:taylor}(b) gives
\[
e^{-x^2}=\sum_{k=0}^m\frac{(-1)^k x^{2k}}{k!}+r_m(x),
\qquad |r_m(x)|\le\frac{x^{2m+2}}{(m+1)!}.
\]
Integrate the finite sum. The error integral is at most $\frac1{(m+1)!(2m+3)}$, and
\[
\left|I-\sum_{k=0}^m\frac{(-1)^k}{k!\,(2k+1)}\right|\le\frac1{(m+1)!(2m+3)}.
\]
At $m=6$ the bound is $\frac1{75600}<10^{-4}$. Nothing here claims that every function equals its Taylor series, and no infinite sum has been passed under the integral sign.

\begin{reminder}[one formula on each side of a singularity]
An antiderivative on an interval has to live on an interval. The function $\frac1x$ is continuous on the disconnected set $\R\setminus\{0\}$. Its antiderivatives are $\ln|x|+C_+$ on the positive half-line and $\ln|x|+C_-$ on the negative half-line, and the two constants are unrelated. In particular, $\int_{-1}^1\frac1x\dd x$ does not converge as an improper integral: the two cutoffs $\int_{-1}^{-\varepsilon}$ and $\int_{\varepsilon}^1$ each diverge as $\varepsilon\to0^+$. If you force the cutoffs to be symmetric, the two pieces cancel and the limit is $0$. That symmetric limit is a principal value. It is a different object from the integral.
\end{reminder}

\photograph{rocket}{Falcon 9 launching DSCOVR, with bright exhaust against a dark sky.}{DSCOVR launch on Falcon 9. SpaceX Photos, public domain; via scikit-image. A launch trajectory needs both an estimate and a bound on the error. The photograph is a reminder of that distinction. The one-variable examples here are not a flight-dynamics model, and the photograph is not an endorsement.}

\section*{What this chapter prepares}
Differentiation and integration undo each other, once the fundamental theorem is in hand. Replacing a function by a polynomial is legitimate when the remainder is bounded, as it is for $\sin(0.1)$ and for the truncated series of $e^{-x^2}$. One recurrence, $W_n=\frac{n-1}n W_{n-2}$, produces Wallis's product, the constant in Stirling's formula, and $\int_0^\infty e^{-x^2}\dd x=\frac{\sqrt\pi}2$.

The next chapter asks a different question. Not what a quantity equals: what it cannot exceed. The number $e$ appears as the best constant in an inequality about averages of a positive sequence. Convexity is the reason, and the same convexity is what singles out $\Gamma$ among all the functions that interpolate $n!$.

Volume~4 turns the Gaussian integral into the law of a sum of many small independent pieces. Volume~5 lives on the formulas proved here: special functions, sharper remainders, Fourier series (Exercise~\ref{ex:rl} is the first step), and the asymptotic series of Proposition~\ref{prop:stirlingnext}. Volume~7 is where the integral in the contract is constructed. Volume~8 takes the transform into infinite-dimensional spaces. The inequalities come first.

\section{Exercises}

\begin{exercise}[$\star$]\label{ex:xnlnx}
Compute $\displaystyle\int_0^1x^n\ln x\dd x$ for $n\in\N$.
\end{exercise}

\begin{exercise}[$\star$]\label{ex:riemann}
Compute the limits of $\displaystyle\sum_{k=1}^n\frac n{n^2+k^2}$ and of $\displaystyle\left(\prod_{k=1}^n\left(1+\frac kn\right)\right)^{1/n}$.
\end{exercise}

\begin{exercise}[$\star\star$]\label{ex:taylorln}
\begin{enumerate}[label=(\alph*)]
\item Show that $e^x=\sum_{k=0}^\infty\frac{x^k}{k!}$ for every real $x$.
\item Integrating $\frac1{1+t}=\sum_{k=0}^{n-1}(-t)^k+\frac{(-t)^n}{1+t}$ over $[0,1]$, give a second proof that $\sum_{k\ge1}\frac{(-1)^{k+1}}k=\ln2$, with an explicit error bound.
\item For $z=x+iy$ with $x,y\in\R$, show that
$\sum_{n\ge0}z^n/n!=e^x(\cos y+i\sin y)$, agreeing with the complex exponential
of Chapter~\ref{ch:complex}. Use the Cauchy product; no complex differentiation is needed.
\end{enumerate}
\end{exercise}

\begin{exercise}[$\star\star$]\label{ex:eirr}
Show that $e$ is irrational.
\end{exercise}

\begin{exercise}[$\star\star$]\label{ex:wallisprod}
\emph{Wallis' product.} Show that $\displaystyle\frac{W_{2p}}{W_{2p+1}}\to1$, and deduce
\[
\begin{aligned}\frac\pi2&=\lim_{p\to\infty}\prod_{k=1}^p\frac{(2k)^2}{(2k-1)(2k+1)}\eqbr=\frac21\cdot\frac23\cdot\frac43\cdot\frac45\cdots\end{aligned}
\]
\end{exercise}

\begin{exercise}[$\star\star$]\label{ex:rl}
\emph{Riemann--Lebesgue.} Let $f$ be of class $C^1$ on $[a,b]$. Show that $\displaystyle\int_a^bf(t)\sin(\lambda t)\dd t\to0$ as $\lambda\to\infty$.
\end{exercise}

\begin{exercise}[$\star\star\star$]\label{ex:gaussian}
\emph{The Gaussian integral, without polar coordinates.}
\begin{enumerate}[label=(\alph*)]
\item Show that $\left(1-\frac{x^2}n\right)^n\le e^{-x^2}$ for $0\le x\le\sqrt n$, and $e^{-x^2}\le\left(1+\frac{x^2}n\right)^{-n}$ for $x\ge0$.
\item Show that $\int_0^{\sqrt n}\left(1-\frac{x^2}n\right)^n\dd x=\sqrt n\,W_{2n+1}$ and $\int_0^\infty\left(1+\frac{x^2}n\right)^{-n}\dd x=\sqrt n\,W_{2n-2}$ (for $n\ge1$).
\item Conclude that $\displaystyle\int_0^\infty e^{-x^2}\dd x=\frac{\sqrt\pi}2$.
\end{enumerate}
\end{exercise}

\subsection*{Further exercises}

\begin{exercise}[$\star$]\label{ex:xnex}
Show that $\displaystyle\int_0^\infty x^ne^{-x}\dd x=n!$ by induction, and that $\displaystyle\int_0^1x^m(\ln x)^n\dd x=\frac{(-1)^nn!}{(m+1)^{n+1}}$ for $m,n\in\N$ (this generalises Exercise~\ref{ex:xnlnx}).
\end{exercise}

\begin{exercise}[$\star$]\label{ex:leibnizpi}
\emph{Leibniz's series (1674).} Integrating $\frac1{1+t^2}=\sum_{k=0}^{n-1}(-1)^kt^{2k}+\frac{(-1)^nt^{2n}}{1+t^2}$ over $[0,x]$, show that $\arctan x=\sum_{k=0}^\infty\frac{(-1)^kx^{2k+1}}{2k+1}$ for $0\le x\le1$, with error at most $\frac{x^{2n+1}}{2n+1}$ after $n$ terms. Deduce $\frac\pi4=1-\frac13+\frac15-\frac17+\cdots$, and explain why this is a terrible way to compute $\pi$ and why $\frac\pi4=4\arctan\frac15-\arctan\frac1{239}$ (Machin, 1706) is a good one.
\end{exercise}

\begin{exercise}[$\star\star$]\label{ex:lhopital}
\begin{enumerate}[label=(\alph*)]
\item Compute $\lim_{x\to0}\frac{x-\sin x}{x^3}$, $\lim_{x\to0}\frac{e^x-1-x}{x^2}$ each by l'H\^opital and by Taylor--Young. Compute also
$\lim_{x\to0^+}x\ln x$ by l'H\^opital and by an elementary exponential bound. Deduce $\lim_{x\to0^+}x^x$.
\item \emph{The $\frac\infty\infty$ form.} Let $f,g$ be differentiable on $]a,a+h[$ with $g'\ne0$, $g(x)\to+\infty$ as $x\to a^+$ and $\frac{f'}{g'}\to\ell\in\R$. Show that $\frac fg\to\ell$. \emph{(Fix $x_0$ close to $a$; for $x<x_0$, Cauchy's mean value theorem on $[x,x_0]$ gives $\frac{f(x)-f(x_0)}{g(x)-g(x_0)}$ close to $\ell$; then let $x\to a$ with $x_0$ fixed.)}
\end{enumerate}
\end{exercise}

\begin{exercise}[$\star\star\star$]\label{ex:trapezoid}
\emph{Bridge: Taylor approximation and roots-of-unity cancellation.}
Apply the trapezoid rule with $n=10$ to $\displaystyle\int_0^1\frac{\dd x}{1+x^2}=\frac\pi4$. Compute the error bound of Proposition~\ref{prop:quadrature} and compare with the actual error ($T_{10}\approx0.784981$, $\frac\pi4\approx0.785398$). Then show that for the \emph{periodic} integral $\int_0^{2\pi}e^{\cos x}\dd x$ the trapezoid rule with $n=10$ is already correct to eight decimals, and find which step of the error estimate becomes an identity for a periodic integrand.
\end{exercise}

\begin{exercise}[$\star\star$]\label{ex:frullani}
\emph{Frullani's integral (1828).} Let $f$ be continuous on $[0,\infty[$ with a limit $f(\infty)$ at infinity, and $a,b>0$. Show that
\[
\int_0^\infty\frac{f(ax)-f(bx)}x\dd x=\bigl(f(0)-f(\infty)\bigr)\ln\frac ba.
\]
Deduce $\int_0^\infty\frac{e^{-ax}-e^{-bx}}x\dd x=\ln\frac ba$ and $\int_0^\infty\frac{\arctan(bx)-\arctan(ax)}{x}\dd x=\frac\pi2\ln\frac ba$.
\end{exercise}

\begin{exercise}[$\star\star$]\label{ex:gammahalf}
\begin{enumerate}[label=(\alph*)]
\item Show that $\displaystyle\int_0^1(1-t^2)^n\dd t=W_{2n+1}$ and $\displaystyle\int_0^1t^{m}(1-t)^{n}\dd t=\frac{m!\,n!}{(m+n+1)!}$ for $m,n\in\N$ (Euler's Beta integral).
\item Deduce that $\displaystyle\int_0^\infty x^{2n}e^{-x^2}\dd x=\frac{(2n)!}{4^nn!}\cdot\frac{\sqrt\pi}2$. These half-line integrals determine the even moments of $e^{-x^2}/\sqrt\pi$, which Volume~4 needs.
\end{enumerate}
\end{exercise}

\begin{exercise}[$\star\star\star$]\label{ex:niven}
\emph{Niven's proof that $\pi$ is irrational (1947).} Suppose $\pi=\frac ab$ with $a,b\in\N^*$, and for $n\ge1$ let $f(x)=\frac{x^n(a-bx)^n}{n!}$ and $F(x)=f(x)-f''(x)+f^{(4)}(x)-\dots+(-1)^nf^{(2n)}(x)$.
\begin{enumerate}[label=(\alph*)]
\item Show that $f^{(k)}(0)$ and $f^{(k)}(\pi)$ are integers for every $k$.
\item Show that $\frac{\dd}{\dd x}\bigl(F'(x)\sin x-F(x)\cos x\bigr)=f(x)\sin x$, and deduce that $I_n=\int_0^\pi f(x)\sin x\dd x$ is an integer.
\item Show that $0<I_n<\frac{\pi(\pi a)^n}{n!}$, and conclude.
\end{enumerate}
\end{exercise}

\begin{exercise}[$\star\star\star$]\label{ex:sophomore}
\emph{The sophomore's dream (Johann Bernoulli, 1697).} Show that
\[
\int_0^1x^{-x}\dd x=\sum_{n=1}^\infty\frac1{n^n}=1.29128\ldots\qquad\text{and}\qquad\int_0^1x^{x}\dd x=\sum_{n=1}^\infty\frac{(-1)^{n+1}}{n^n}.
\]
\emph{(Write $x^{-x}=e^{-x\ln x}$, expand the exponential with the explicit remainder of Exercise~\ref{ex:taylorln}(a), and use Exercise~\ref{ex:xnex}.)}
\end{exercise}

\begin{challenge}
\begin{enumerate}[label=(\arabic*)]
\item \emph{The Basel problem by a double integral.} Show that $\zeta(2)=\int_0^1\!\!\int_0^1\frac{\dd x\dd y}{1-xy}$, and evaluate the double integral by the change of variables $x=\frac{u-v}{\sqrt2}$, $y=\frac{u+v}{\sqrt2}$ (rotate the square by $45^\circ$): the integrand becomes $\frac2{2-u^2+v^2}$, whose integral in $v$ is an $\arctan$. This is Apostol's 1983 proof. Volume~7 supplies the two theorems on double integrals it needs.
\item \emph{$\pi$ is transcendental.} Lindemann proved in 1882 that $\pi$ is not a root of any polynomial with integer coefficients, which settled the ancient question of squaring the circle. His proof refines Hermite's 1873 proof for $e$, which refines the integration-by-parts trick of Exercises~\ref{ex:eirr} and~\ref{ex:niven}. Read Hermite's proof for $e$ (in Hardy and Wright, Chapter~XI) and locate the trick.
\item \emph{Open.} Nobody knows whether $e+\pi$, $e\pi$ or $\pi^e$ is irrational. Show that if $e+\pi$ and $e\pi$ were both rational, then $e$ and $\pi$ would be roots of one quadratic with rational coefficients, hence algebraic. Either conclusion contradicts a transcendence theorem: Hermite proved in 1873 that $e$ is transcendental, and Lindemann proved in 1882 that $\pi$ is. Those proofs are outlined in the notes; they are not proved in this chapter. The parenthetical quadratic step is the part to write out.
\end{enumerate}
Notes at the end of the chapter.
\end{challenge}

\begin{beauty}
One recurrence, $W_n=\frac{n-1}nW_{n-2}$, gives Wallis' product, the constant in Stirling's formula, and the Gaussian integral. The three share that relation. In Volume~4 they meet again in the central limit theorem: the $\sqrt{2\pi}$ in the normal density is Stirling's $\sqrt{2\pi}$. Euler built $\Gamma$ to interpolate $n!$. Volume~5 puts that same function next to $\zeta(s)$ in Riemann's functional equation. Riemann--Lebesgue (Exercise~\ref{ex:rl}) is the first step of the Fourier analysis in Volume~5. Niven's argument (Exercise~\ref{ex:niven}) is a one-page proof that $\pi$ is irrational, a fact Lambert reached in 1768 by a much longer route.
\end{beauty}

\begin{exercise}[$\star\star$]\label{pb:lightsout}
Lamps numbered by $\Z/n\Z$, $n\ge3$, stand on a circle, each on or off. Pressing lamp $i$ switches the state of the lamps
$i-1$, $i$ and $i+1$.
\begin{enumerate}[label=(\alph*)]
\item Show that every configuration can be switched off if and only if $3\nmid n$.
\item Let $3\mid n$. For $r\in\{0,1,2\}$, let $S_r$ be the number of lit lamps $i$ with $i\equiv r\pmod3$. Show that a
configuration can be switched off if and only if $S_0$, $S_1$ and $S_2$ have the same parity. What proportion of the
configurations is that?
\end{enumerate}
\end{exercise}

\begin{exercise}[$\star\star$]\label{pb:hyperplanes}
\emph{Covering a cube minus a point.}
\begin{enumerate}[label=(\alph*)]
\item (Alon and Füredi, 1993.) Show that affine hyperplanes of $\mathbb F_q^n$ that cover every point except $0$ must number at
least $n(q-1)$, and that $n(q-1)$ is enough.
\item (Jamison 1977, Brouwer and Schrijver 1978.) A \emph{blocking set} is a set of points of $\mathbb F_q^n$ that meets every affine
hyperplane. Show that a blocking set has at least $n(q-1)+1$ points. Find one of that size.
\end{enumerate}
\end{exercise}

\begin{exercise}[$\star\star$]\label{pb:kakeya}
\emph{Dvir's Kakeya theorem.} A \emph{Kakeya set} $K\subset\mathbb F_q^n$ contains a line in every direction: for every
$v\ne0$ there is $a$ with $\{a+tv:t\in\mathbb F_q\}\subset K$.
\begin{enumerate}[label=(\alph*)]
\item (Dvir, 2008.) Show that $|K|\ge\binom{q+n-1}n$. In particular $|K|\ge\frac{q^n}{n!}$: a Kakeya set fills a fixed fraction
of the space.
\item For $q$ odd, build a Kakeya set in $\mathbb F_q^2$ with $\frac{q(q+1)}2+\frac{q-1}2$ points.
\end{enumerate}
\end{exercise}

\section{Hints}

\begin{hint}{ex:xnlnx}
Integrate by parts on $[\varepsilon,1]$, differentiating $\ln x$.
\end{hint}
\begin{hint}{ex:riemann}
$\frac n{n^2+k^2}=\frac1n\cdot\frac1{1+(k/n)^2}$. For the product, take logarithms.
\end{hint}
\begin{hint}{ex:taylorln}
(a) Taylor--Lagrange with $\sup|\exp^{(n+1)}|\le e^{|x|}$ on $[0,x]$. (b) The error term is $\int_0^1\frac{t^n}{1+t}\dd t$; bound $\frac1{1+t}$ by $1$. (c) Apply Taylor to sine and cosine, split the series at $iy$
into even and odd terms, and prove the series addition law with the Cauchy product.
\end{hint}
\begin{hint}{ex:eirr}
If $e=\frac pq$, show that $q!\left(e-\sum_{k=0}^q\frac1{k!}\right)$ is an integer strictly between $0$ and $1$.
\end{hint}
\begin{hint}{ex:wallisprod}
$1\le\frac{W_{2p}}{W_{2p+1}}\le\frac{W_{2p-1}}{W_{2p+1}}$, and the last ratio is given by the recurrence.
\end{hint}
\begin{hint}{ex:rl}
Integrate by parts: $\sin(\lambda t)$ has the small antiderivative $-\frac{\cos(\lambda t)}\lambda$.
\end{hint}
\begin{hint}{ex:gaussian}
(a) $\ln(1+u)\le u$. (b) Substitute $x=\sqrt n\sin\theta$, then $x=\sqrt n\tan\theta$. (c) Squeeze, and use $W_n\sim\sqrt{\frac\pi{2n}}$.
\end{hint}
\begin{hint}{ex:xnex}
Parts, differentiating the power; for the second integral, parts differentiating $(\ln x)^n$, on $[\varepsilon,1]$.
\end{hint}
\begin{hint}{ex:leibnizpi}
The remainder is $\int_0^x\frac{t^{2n}}{1+t^2}\dd t\le\int_0^xt^{2n}\dd t$. For Machin, compute $\tan$ of $4\arctan\frac15$ with the addition formula, twice.
\end{hint}
\begin{hint}{ex:lhopital}
(a) For $x^x$, write $e^{x\ln x}$. (b) Write $\frac{f(x)}{g(x)}=\frac{f(x)-f(x_0)}{g(x)-g(x_0)}\cdot\frac{g(x)-g(x_0)}{g(x)}+\frac{f(x_0)}{g(x)}$, and note that the middle factor and the last term tend to $1$ and $0$ as $x\to a^+$.
\end{hint}
\begin{hint}{ex:trapezoid}
$f''(x)=\frac{6x^2-2}{(1+x^2)^3}$, bounded by $2$ on $[0,1]$. For the periodic case, identify the vanishing endpoint difference in the
kernel proof. To certify the stated precision, expand $e^u$ through degree $19$
at $u=\cos x$ and use the ten-point roots-of-unity filter on each power of cosine.
\end{hint}
\begin{hint}{ex:frullani}
On $[\varepsilon,A]$, substitute in each integral separately: $\int_\varepsilon^A\frac{f(ax)}x\dd x=\int_{a\varepsilon}^{aA}\frac{f(u)}u\dd u$. The difference reduces to $\int_{a\varepsilon}^{b\varepsilon}\frac{f(u)}u\dd u-\int_{aA}^{bA}\frac{f(u)}u\dd u$, where $f$ is nearly constant.
\end{hint}
\begin{hint}{ex:gammahalf}
(a) $t=\sin\theta$ for the first; for the second, integrate by parts to get a recurrence in $n$, or substitute $t=\sin^2\theta$ and use Wallis. (b) $x^2=t$ turns the integral into $\frac12\Gamma\left(n+\frac12\right)$.
\end{hint}
\begin{hint}{ex:niven}
(a) $f$ has integer coefficients divided by $n!$ and only powers $x^n,\dots,x^{2n}$; and $f(\pi-x)=f(x)$. (b) Differentiate, and note $F+F''=f$ because $f^{(2n+2)}=0$. (c) On $]0,\pi[$, $0<x(a-bx)<\pi a$ and $0<\sin x\le1$; then $\frac{c^n}{n!}\to0$.
\end{hint}
\begin{hint}{ex:sophomore}
$|x\ln x|\le\frac1e$ on $[0,1]$, so the Taylor--Lagrange remainder of the exponential after $N$ terms is at most $\frac{e^{1/e}}{(N+1)!\,e^{N+1}}$, uniformly in $x$; integrate term by term with $\int_0^1\frac{(-x\ln x)^n}{n!}\dd x=\frac1{(n+1)^{n+1}}$.
\end{hint}

\begin{hint}{pb:lightsout}
Identify configurations with $\mathbb F_2[x]/(x^n-1)$: pressing the set $u$ produces $(x^{-1}+1+x)u$. For (b), the matrix is
symmetric, so its image is the orthogonal of its kernel.
\end{hint}

\begin{hint}{pb:hyperplanes}
(a) Compare $\prod_i(1-a_i\cdot x)$ with $\prod_j\prod_{s\ne0}(x_j-s)$. (b) The hyperplanes $\{x:a\cdot x=1\}$ avoid $0$; swap the
roles of $a$ and $x$.
\end{hint}

\begin{hint}{pb:kakeya}
(a) A non-zero polynomial of degree at most $q-1$ vanishing on $K$ exists if $K$ is small. Restrict it to a line and look at the
top coefficient. (b) Use the lines $y=mx+m^2$.
\end{hint}

\section{Solutions}
\begin{solution}{ex:xnlnx}
\idea A product of a power and a logarithm: integrate by parts so as to \emph{differentiate} the logarithm, which turns it into $\frac1x$ and leaves a pure power. The integral is improper at $0$, so work on $[\varepsilon,1]$ and let $\varepsilon\to0$.
\writeup On $[\varepsilon,1]$, integrate by parts with $u'=x^n$, $v=\ln x$:
$\int_\varepsilon^1x^n\ln x\dd x=\left[\frac{x^{n+1}}{n+1}\ln x\right]_\varepsilon^1-\int_\varepsilon^1\frac{x^n}{n+1}\dd x$.
As $\varepsilon\to0$, $\varepsilon^{n+1}\ln\varepsilon\to0$, and the integral equals $-\dfrac1{(n+1)^2}$.
\end{solution}

\begin{solution}{ex:riemann}
\idea A sum of $n$ terms whose general term is a function of $\frac kn$, times $\frac1n$, is a Riemann sum. Factor out $n^2$ in the first to see $\frac1n\cdot\frac1{1+(k/n)^2}$. For a product, take logarithms to get a sum.
\writeup
\[
\sum_{k=1}^n\frac n{n^2+k^2}=\frac1n\sum_{k=1}^n\frac1{1+(k/n)^2}\to\int_0^1\frac{\mathrm dx}{1+x^2}=\frac\pi4
\]
 by Proposition~\ref{prop:inttools}(d).
For the product, take logarithms:
\[
\frac1n\sum_{k=1}^n\ln\left(1+\frac kn\right)\to\int_0^1\ln(1+x)\dd x=\left[(1+x)\ln(1+x)-x\right]_0^1=2\ln2-1,
\]
so the limit is $e^{2\ln2-1}=\frac4e$ by continuity of $\exp$.
\end{solution}

\begin{solution}{ex:taylorln}
\idea (a) Showing that a function equals its Taylor series means showing the remainder tends to $0$; for $\exp$ all derivatives are $\exp$, so Taylor--Lagrange gives a bound $\frac{|x|^{n+1}}{(n+1)!}e^{|x|}$, and factorials beat powers. (This reconciles the real exponential of analysis, the solution of $y'=y$, $y(0)=1$, with the series representation of $e^{i\theta}$, identified in part~(c) below.) (b) The identity is a finite geometric sum, so it is exact; integrating it gives the partial sum of the alternating series plus an explicit error integral, easy to bound.
\writeup (a) Apply Theorem~\ref{thm:taylor}(b) at $a=0$. Every derivative of $\exp$ is $\exp$, and $\exp\le e^{|x|}$ between $0$ and $x$, so the error is at most $\frac{|x|^{n+1}}{(n+1)!}e^{|x|}$, which tends to $0$ because the ratio of consecutive terms, $\frac{|x|}{n+2}$, tends to $0$.

(b) Integrating over $[0,1]$: $\ln2=\sum_{k=0}^{n-1}\frac{(-1)^k}{k+1}+(-1)^n\int_0^1\frac{t^n}{1+t}\dd t$, and
$0\le\int_0^1\frac{t^n}{1+t}\dd t\le\int_0^1t^n\dd t=\frac1{n+1}$. So the partial sums of $\sum\frac{(-1)^{k+1}}k$ are within $\frac1{n+1}$ of $\ln2$.

(c) Put $E(z)=\sum_{n\ge0}z^n/n!$. At $z=0$ its sum is $1$; for $z\ne0$ the ratio test proves absolute convergence.
Taylor--Lagrange for sine and cosine bounds each remainder at a real $y$ by
$|y|^{N+1}/(N+1)!$, tending to zero. Their Taylor coefficients therefore give
\[
E(iy)=\sum_{k\ge0}\frac{(-1)^k y^{2k}}{(2k)!}
+i\sum_{k\ge0}\frac{(-1)^k y^{2k+1}}{(2k+1)!}
=\cos y+i\sin y.
\]
Separating the even and odd terms is justified by absolute convergence.
The complex Cauchy product and the binomial formula give
\[
E(a)E(b)=\sum_{n\ge0}\sum_{k=0}^n
\frac{a^kb^{n-k}}{k!(n-k)!}
=\sum_{n\ge0}\frac{(a+b)^n}{n!}=E(a+b).
\]
Part (a) gives $E(x)=e^x$ for real $x$. Hence
$E(x+iy)=E(x)E(iy)=e^x(\cos y+i\sin y)$, as required.
\end{solution}

\begin{solution}{ex:eirr}
\idea Irrationality proofs follow one pattern: assume $x=\frac pq$, build an integer from it, and show that integer lies strictly between $0$ and $1$. The series of $e$ converges very fast, so the tail after $\frac1{q!}$ is tiny; multiplying by $q!$ makes the head an integer and keeps the tail below $1$.
\writeup Suppose $e=\frac pq$ with $p,q\ge1$. Then $A=q!\left(e-\sum_{k=0}^q\frac1{k!}\right)$ is an integer, since $q!e=(q-1)!\,p$ and $\frac{q!}{k!}\in\N$ for $k\le q$. But
\[
\begin{aligned}0<A&=\sum_{k>q}\frac{q!}{k!}=\frac1{q+1}+\frac1{(q+1)(q+2)}+\cdots\eqbr<\sum_{j\ge1}\frac1{(q+1)^j}=\frac1q\le1.\end{aligned}
\]
No integer lies strictly between $0$ and $1$.
\end{solution}

\begin{solution}{ex:wallisprod}
\idea The ratio $\frac{W_{2p}}{W_{2p+1}}$ is squeezed between $1$ and $\frac{W_{2p-1}}{W_{2p+1}}$ by monotonicity, and the recurrence computes the latter exactly. Then write the ratio with the closed forms (ii) and reorganise the factorials into a product over $k$.
\writeup By monotonicity, $W_{2p+1}\le W_{2p}\le W_{2p-1}$, so $1\le\frac{W_{2p}}{W_{2p+1}}\le\frac{W_{2p-1}}{W_{2p+1}}=\frac{2p+1}{2p}\to1$. By Proposition~\ref{prop:wallis}(ii),
\[
\begin{aligned}\frac{W_{2p}}{W_{2p+1}}&=\frac\pi2\cdot\frac{(2p)!\,(2p+1)!}{16^p\,(p!)^4}\eqbr=\frac\pi2\prod_{k=1}^p\frac{(2k-1)(2k+1)}{(2k)^2},\end{aligned}
\]
because $(2p)!=\prod_{k=1}^p(2k-1)(2k)$, $(2p+1)!=(2p)!\,(2p+1)$ and $\prod_{k=1}^p(2k)=2^pp!$. Letting $p\to\infty$ gives the product formula.
\end{solution}

\begin{solution}{ex:rl}
\idea A rapidly oscillating factor $\sin\lambda t$ has a small antiderivative, $-\frac{\cos\lambda t}\lambda$. Integrating by parts moves the derivative onto $f$ (which is why we assume $C^1$) and produces a factor $\frac1\lambda$.
\writeup Integrate by parts:
\[\begin{aligned}\int_a^bf(t)\sin(\lambda t)\dd t&=\left[-f(t)\frac{\cos\lambda t}\lambda\right]_a^b\eqbr+\frac1\lambda\int_a^bf'(t)\cos(\lambda t)\dd t.\end{aligned}\]
Both terms together are at most $\frac C\lambda$ in absolute value, where $C=|f(a)|+|f(b)|+\int_a^b|f'|$, so they tend to $0$.
\end{solution}

\begin{solution}{ex:gaussian}
\idea We know $\left(1\pm\frac{x^2}n\right)^{\pm n}\to e^{-x^2}$, and the integrals of these approximations are Wallis integrals after a trigonometric substitution ($1-\sin^2=\cos^2$ and $1+\tan^2=\frac1{\cos^2}$). So squeeze $e^{-x^2}$ between them, and let the known asymptotics $W_n\sim\sqrt{\frac\pi{2n}}$ finish. The inequalities in (a) are the concavity bound $\ln(1+u)\le u$.
\writeup (a) With $u=\frac{x^2}n$: $\ln(1-u)\le-u$ for $0\le u<1$ gives $n\ln(1-u)\le-x^2$, the first inequality (it is trivial at $u=1$); $\ln(1+u)\le u$ gives $(1+u)^n\le e^{nu}=e^{x^2}$, the second.

(b) Substitute $x=\sqrt n\sin\theta$, $\theta\in\left[0,\frac\pi2\right]$: $\left(1-\frac{x^2}n\right)^n\dd x=\cos^{2n}\theta\cdot\sqrt n\cos\theta\dd\theta$, giving $\sqrt nW_{2n+1}$.
Substitute $x=\sqrt n\tan\theta$, $\theta\in\left[0,\frac\pi2\right[$: $\left(1+\frac{x^2}n\right)^{-n}\dd x=\cos^{2n}\theta\cdot\frac{\sqrt n}{\cos^2\theta}\dd\theta$, giving $\sqrt nW_{2n-2}$ (an improper integral; apply the change of variables on $[0,A]$ and let $A\to\infty$).

(c) Let $I=\int_0^\infty e^{-x^2}\dd x$ (it converges since $e^{-x^2}\le e^{-x}$ for $x\ge1$). By (a) and (b),
$\sqrt nW_{2n+1}\le\int_0^{\sqrt n}e^{-x^2}\dd x\le I\le\sqrt nW_{2n-2}$. By Proposition~\ref{prop:wallis}(iii),
$\sqrt nW_{2n+1}\sim\sqrt n\sqrt{\frac\pi{4n}}=\frac{\sqrt\pi}2$, and likewise for $\sqrt nW_{2n-2}$. Hence $I=\frac{\sqrt\pi}2$.
\end{solution}


\begin{solution}{ex:xnex}
\idea Both are the integration-by-parts machine: each step lowers the exponent of the power (first integral) or of the logarithm (second) by one, so the answers are products of the successive exponents, that is, factorials. The improper ends ($\infty$ and $0$) are harmless because $x^ne^{-x}$ and $x^{m+1}(\ln x)^n$ tend to $0$ there.
\writeup Let $I_n=\int_0^\infty x^ne^{-x}\dd x$. Then $I_0=1$, and for $n\ge1$, on $[0,A]$, with $u'=e^{-x}$, $v=x^n$: $\int_0^Ax^ne^{-x}\dd x=\left[-x^ne^{-x}\right]_0^A+n\int_0^Ax^{n-1}e^{-x}\dd x$; as $A\to\infty$, $A^ne^{-A}\to0$, so $I_n=nI_{n-1}$ and $I_n=n!$.
Let $J_{m,n}=\int_0^1x^m(\ln x)^n\dd x$. Then $J_{m,0}=\frac1{m+1}$, and for $n\ge1$, on $[\varepsilon,1]$ with $u'=x^m$, $v=(\ln x)^n$: $\int_\varepsilon^1x^m(\ln x)^n\dd x=\left[\frac{x^{m+1}}{m+1}(\ln x)^n\right]_\varepsilon^1-\frac n{m+1}\int_\varepsilon^1x^m(\ln x)^{n-1}\dd x$. The bracket tends to $0$ (as $\varepsilon^{m+1}(\ln\varepsilon)^n\to0$), so $J_{m,n}=-\frac n{m+1}J_{m,n-1}$, and $J_{m,n}=\frac{(-1)^nn!}{(m+1)^{n+1}}$.
\end{solution}

\begin{solution}{ex:leibnizpi}
\idea Same recipe as for $\ln2$: a finite geometric sum with an explicit remainder, integrated. The remainder is tiny for small $x$ and merely small for $x=1$, which is why Leibniz's series is slow and Machin's fast: $\frac15$ and $\frac1{239}$ are small.
\writeup Integrating the identity from $0$ to $x$ and using $\arctan'=\frac1{1+t^2}$ (fundamental theorem): $\arctan x=\sum_{k=0}^{n-1}\frac{(-1)^kx^{2k+1}}{2k+1}+(-1)^nR_n(x)$ with $0\le R_n(x)=\int_0^x\frac{t^{2n}}{1+t^2}\dd t\le\int_0^xt^{2n}\dd t=\frac{x^{2n+1}}{2n+1}$. For $0\le x\le1$ this tends to $0$, which proves the series and the error bound. At $x=1$: $\frac\pi4=1-\frac13+\frac15-\cdots$, James Gregory's and Leibniz's series (1671, 1674), known to Madhava in Kerala around 1400.
For $\pi/4$ the error after $n$ terms is at most $1/(2n+1)$; for $\pi$ the bound is $4/(2n+1)$. To guarantee absolute error below $10^{-10}$ for $\pi$ by this bound requires $n\ge2\cdot10^{10}$. With Machin: $\tan\left(2\arctan\frac15\right)=\frac{2/5}{1-1/25}=\frac5{12}$, $\tan\left(4\arctan\frac15\right)=\frac{10/12}{1-25/144}=\frac{120}{119}$, and $\tan\left(4\arctan\frac15-\frac\pi4\right)=\frac{120/119-1}{1+120/119}=\frac1{239}$; since $4\arctan\frac15-\frac\pi4\in\,]0,\frac\pi2[$, this gives Machin's formula. The series for $\arctan\frac15$ has error at most $\frac{5^{-(2n+1)}}{2n+1}$ after $n$ terms: $7$ terms of it and $2$ of $\arctan\frac1{239}$ give a combined error in $\pi$ below $10^{-10}$ after the factors $16$ and $4$ are included. Machin computed $100$ decimals this way in 1706, and formulas of the same kind held the record for more than 270 years: the first computer calculation, $2037$ decimals on the ENIAC in 1949, used Machin's own formula, and only in 1982 (Tamura and Kanada, with Gauss's arithmetic-geometric mean) did the record pass to other methods.
\end{solution}

\begin{solution}{ex:lhopital}
\idea (a) is routine with either tool; the point is to see that Taylor--Young does the first two limits in two lines where l'H\^opital needs several differentiations each. (b) The $\frac\infty\infty$ form cannot use the trick of extending by $0$ at $a$; instead one compares $f$ and $g$ to their values at a fixed nearby point $x_0$, where Cauchy's theorem controls the increments, and then lets $g(x)\to\infty$ wash out the fixed quantities.
\writeup (a) By Taylor--Young at $0$: $x-\sin x=\frac{x^3}6+o(x^3)$, so the first limit is $\frac16$; $e^x-1-x=\frac{x^2}2+o(x^2)$, so the second is $\frac12$. By l'H\^opital, the first needs three differentiations ($\frac{1-\cos x}{3x^2}$, $\frac{\sin x}{6x}$, $\frac{\cos x}6$) and the second two. For $x\ln x=\frac{\ln x}{1/x}$, the $\frac\infty\infty$ form of (b) applies (with $x\to0^+$, $g=\frac1x\to+\infty$): $\frac{1/x}{-1/x^2}=-x\to0$. For an independent bound, put $t=-\ln x\to\infty$. Taylor's theorem gives
$e^t\ge t^2/2$ for $t>0$, so $|x\ln x|=te^{-t}\le2/t\to0$.
Hence $x^x=e^{x\ln x}\to e^0=1$.

(b) Let $\varepsilon>0$ and choose $x_0\in\,]a,a+h[$ with $\left|\frac{f'(t)}{g'(t)}-\ell\right|\le\varepsilon$ for $a<t<x_0$. For $a<x<x_0$, Cauchy's mean value theorem on $[x,x_0]$ gives $t\in\,]x,x_0[$ with $\frac{f(x)-f(x_0)}{g(x)-g(x_0)}=\frac{f'(t)}{g'(t)}$ (note $g(x)\ne g(x_0)$ by Rolle, as $g'\ne0$), so this quotient is within $\varepsilon$ of $\ell$. Now
\[
\frac{f(x)}{g(x)}=\frac{f(x)-f(x_0)}{g(x)-g(x_0)}\cdot\left(1-\frac{g(x_0)}{g(x)}\right)+\frac{f(x_0)}{g(x)}.
\]
As $x\to a^+$, $g(x)\to\infty$, so the factor in parentheses tends to $1$ and the last term to $0$, while the first quotient stays within $\varepsilon$ of $\ell$. Hence $\left|\frac{f(x)}{g(x)}-\ell\right|\le2\varepsilon$ for $x$ close enough to $a$.
\end{solution}

\begin{solution}{ex:trapezoid}
\idea The error bound is a worst case; the actual error is smaller because $f''$ changes sign on $[0,1]$. For a periodic integrand, extracting the mean of the error kernel produces a leading term proportional to $f'(b)-f'(a)=0$. For a certified ten-point estimate, a different language is more efficient:
Taylor approximation plus the roots-of-unity filter makes the error nearly vanish
term by term. This combines complex algebra with quantitative calculus.
\writeup With $f(x)=\frac1{1+x^2}$: $f'(x)=\frac{-2x}{(1+x^2)^2}$ and $f''(x)=\frac{6x^2-2}{(1+x^2)^3}$, which lies between $-2$ (at $0$) and $\frac12$ (at $1$), so $M=2$. The bound is $\frac{2\cdot1}{12\cdot100}=\frac1{600}\approx0.00167$; the actual error is $\frac\pi4-T_{10}\approx0.000417$, four times smaller, because $f''$ is negative near $0$ and positive near $1$, and the two parts of the error partly cancel.

For the periodic case, first identify the cancellation in the error kernel.
On each subinterval write $p=h^2/12+\widetilde p$, where
$\int\widetilde p=0$. The constant parts of $-\int f''p$ sum to
$-h^2(f'(2\pi)-f'(0))/12=0$ for $f(x)=e^{\cos x}$.
This explains an improvement, but does not yet certify eight decimals.
Roots of unity give a sharper bound.

\emph{How to find the certificate.} The ten equally spaced points annihilate
oscillations $e^{imx}$ unless $10$ divides $m$. Taylor turns $e^{\cos x}$ into
a finite sum of such oscillations, so only a few terms can contribute to the error.

\emph{Write-up.} Let $P(u)=\sum_{r=0}^{19}u^r/r!$. For $-1\le u\le1$,
Taylor--Lagrange and $e<3$ give $|e^u-P(u)|<3/20!$.
Write $I(g)=\int_0^{2\pi}g(x)\,\dd x$ and
$T(g)=(2\pi/10)\sum_{k=0}^9g(2\pi k/10)$; periodicity makes this the
trapezoid rule. The approximation changes $T(g)-I(g)$ by less than $12\pi/20!$.
By the binomial formula,
\[
\cos^r x=2^{-r}\sum_{j=0}^r\binom rj e^{i(r-2j)x}.
\]
The integral of $e^{imx}$ is zero for nonzero integer $m$, by integrating
its real and imaginary parts. The ten-point average is zero unless $10\mid m$,
by the roots-of-unity filter. Since $r\le19$, only frequencies $10$ and $-10$
remain in the difference, giving
\[
T(P\circ\cos)-I(P\circ\cos)
=4\pi\sum_{j=0}^4\frac1{2^{10+2j}j!(10+j)!}.
\]
Now $(10+j)!\ge10!\,11^j$, hence this positive sum is at most
\[
\frac{4\pi}{2^{10}10!}\sum_{j=0}^4\frac1{44^j j!}
<\frac{16}{2^{10}10!}\frac{44}{43},
\]
using $\pi<4$ and a geometric-series bound. Here $\pi<4$ follows already from
$\pi/4=\int_0^1(1+x^2)^{-1}\,\dd x<1$.
Together with the uniform remainder,
\[
|T(e^{\cos x})-I(e^{\cos x})|
<\frac{16\cdot44}{43\cdot2^{10}10!}+\frac{48}{20!}
<4.42\times10^{-9}<\frac12\,10^{-8}.
\]
This is the promised eight-decimal absolute error guarantee. Numerically,
$T_{10}=7.954926524\ldots$ and $I=7.954926521\ldots$; the certificate, not the
agreement of two decimals, establishes the accuracy. No infinite sum has been
interchanged with an integral.
\end{solution}

\begin{solution}{ex:frullani}
\idea The two integrals $\int\frac{f(ax)}x$ and $\int\frac{f(bx)}x$ diverge separately (at $0$ if $f(0)\ne0$, at $\infty$ if $f(\infty)\ne0$), but after the substitution $u=ax$ and $u=bx$ they become the same integral $\int\frac{f(u)}u\dd u$ over slightly different ranges, and the difference is two short integrals, one near $0$ and one near $\infty$, where $f$ is nearly constant and $\int\frac{\dd u}u=\ln\frac ba$.
\writeup For $0<\varepsilon<A$, substituting $u=ax$ and $u=bx$ (Proposition~\ref{prop:inttools}(c)),
\[
\int_\varepsilon^A\frac{f(ax)-f(bx)}x\dd x=\int_{a\varepsilon}^{aA}\frac{f(u)}u\dd u-\int_{b\varepsilon}^{bA}\frac{f(u)}u\dd u=\int_{a\varepsilon}^{b\varepsilon}\frac{f(u)}u\dd u-\int_{aA}^{bA}\frac{f(u)}u\dd u
\]
(assume $a<b$; the two integrals over $[b\varepsilon,aA]$ cancel). On $[a\varepsilon,b\varepsilon]$, $f(u)=f(0)+o(1)$ uniformly as $\varepsilon\to0$ (continuity at $0$), so the first integral is $\bigl(f(0)+o(1)\bigr)\ln\frac ba$; likewise the second is $\bigl(f(\infty)+o(1)\bigr)\ln\frac ba$ as $A\to\infty$. Letting $\varepsilon\to0$ and $A\to\infty$ gives the formula, and in particular shows that the improper integral converges.
With $f(x)=e^{-x}$: $f(0)-f(\infty)=1$, giving $\ln\frac ba$. With $f(x)=-\arctan x$: $f(0)-f(\infty)=\frac\pi2$, giving $\frac\pi2\ln\frac ba$.
\end{solution}

\begin{solution}{ex:gammahalf}
\idea These integrals all reduce to Wallis integrals by the substitutions $t=\sin\theta$ or $t=\sin^2\theta$, which turn $1-t^2$ or $1-t$ into a $\cos^2$. The Beta integral is best found by a recurrence in $n$ obtained by parts, exactly as for $W_n$. The Gaussian moments come from $\Gamma$ by $x^2=t$.
\writeup (a) With $t=\sin\theta$, $\dd t=\cos\theta\dd\theta$: $\int_0^1(1-t^2)^n\dd t=\int_0^{\pi/2}\cos^{2n+1}\theta\dd\theta=W_{2n+1}$.
Let $B(m,n)=\int_0^1t^m(1-t)^n\dd t$. Then $B(m,0)=\frac1{m+1}$, and for $n\ge1$, integrating by parts with $u'=t^m$, $v=(1-t)^n$: $B(m,n)=\left[\frac{t^{m+1}}{m+1}(1-t)^n\right]_0^1+\frac n{m+1}B(m+1,n-1)=\frac n{m+1}B(m+1,n-1)$. Iterating, $B(m,n)=\frac{n}{m+1}\cdot\frac{n-1}{m+2}\cdots\frac1{m+n}\cdot B(m+n,0)=\frac{n!\,m!}{(m+n)!}\cdot\frac1{m+n+1}=\frac{m!\,n!}{(m+n+1)!}$.

(b) With $t=x^2$, $\dd t=2x\dd x$: $\int_0^\infty x^{2n}e^{-x^2}\dd x=\frac12\int_0^\infty t^{n-1/2}e^{-t}\dd t=\frac12\Gamma\left(n+\frac12\right)=\frac{(2n)!}{4^nn!}\cdot\frac{\sqrt\pi}2$ by Proposition~\ref{prop:gammafn}(b). For $n=0$ this is the Gaussian integral; for $n=1$, $\frac{\sqrt\pi}4$. The density $e^{-x^2}/\sqrt\pi$ on $\R$ has total mass $1$, so its even moments are
\[
\frac{2}{\sqrt\pi}\int_0^\infty x^{2n}e^{-x^2}\dd x=\frac{(2n)!}{4^nn!}.
\]
The product $1\cdot3\cdot5\cdots(2n-1)=\frac{(2n)!}{2^nn!}$ is the even moment of the standard normal density $(2\pi)^{-1/2}e^{-x^2/2}$. It follows from the integral above by the substitution $x=u/\sqrt2$, not by dividing by the mass of $e^{-x^2}$.
\end{solution}

\begin{solution}{ex:niven}
\idea The pattern of Exercise~\ref{ex:eirr} once more: from the assumption $\pi=\frac ab$ build an integer, here $I_n$, and show it lies strictly between $0$ and $1$ for $n$ large. The polynomial $f$ is designed so that its derivatives at $0$ and at $\pi$ are integers (the $\frac1{n!}$ is exactly cancelled by the $k!$ coming out of differentiation), and the alternating sum $F$ is designed so that $F''+F=f$, which makes $\int f\sin x$ computable by parts in one stroke.
\writeup (a) Expand $f(x)=\frac1{n!}\sum_{j=n}^{2n}c_jx^j$ with integers $c_j$ (the coefficients of $x^n(a-bx)^n$). Then $f^{(k)}(0)=\frac{k!\,c_k}{n!}$ for $n\le k\le2n$ and $0$ otherwise; since $k\ge n$, $\frac{k!}{n!}$ is an integer, and so is $f^{(k)}(0)$. Next, $f(\pi-x)=\frac{(\pi-x)^n(a-b\pi+bx)^n}{n!}=\frac{(\pi-x)^n(bx)^n}{n!}=f(x)$ because $a=b\pi$; differentiating $k$ times, $f^{(k)}(\pi-x)=(-1)^kf^{(k)}(x)$, so $f^{(k)}(\pi)=(-1)^kf^{(k)}(0)$ is an integer.

(b) $\frac{\dd}{\dd x}\bigl(F'\sin x-F\cos x\bigr)=F''\sin x+F'\cos x-F'\cos x+F\sin x=(F''+F)\sin x$, and $F''+F=f$ because the series defining $F$ telescopes under $F''$: $F''=f''-f^{(4)}+\dots+(-1)^nf^{(2n+2)}$ and $f^{(2n+2)}=0$ ($\deg f=2n$). Hence
$I_n=\bigl[F'(x)\sin x-F(x)\cos x\bigr]_0^\pi=F(\pi)+F(0)$, an integer by (a).

(c) On $]0,\pi[$, $x>0$, $a-bx=b(\pi-x)>0$ and $\sin x>0$, so $f(x)\sin x>0$ and $I_n>0$. Also $x(a-bx)\le\pi a$ (as $x\le\pi$ and $a-bx\le a$) and $\sin x\le1$, so $I_n\le\pi\frac{(\pi a)^n}{n!}$. Since $\frac{c^n}{n!}\to0$ for every $c$ (ratio test), $I_n<1$ for $n$ large: an integer strictly between $0$ and $1$. So $\pi$ is irrational.
Niven's note in the \emph{Bulletin of the AMS} of 1947 is one page long. Lambert had proved the theorem in 1768, by continued fractions, in a much longer argument.
\end{solution}

\begin{solution}{ex:sophomore}
\idea Expand $x^{-x}=e^{-x\ln x}$ as a series in powers of $x\ln x$ and integrate term by term; each term is the integral of Exercise~\ref{ex:xnex}, and the sum collapses to $\sum\frac1{n^n}$. The interchange of sum and integral is justified exactly as in Exercise~\ref{ex:eulerbasel}: the remainder of the exponential series is bounded uniformly on $[0,1]$, because $|x\ln x|\le\frac1e$ there (the maximum of $-x\ln x$ is at $x=\frac1e$).
\writeup Put $u=-x\ln x\in\left[0,\frac1e\right]$ on $]0,1]$ (extend continuously by $0$ at $x=0$). Taylor--Lagrange for $\exp$ at $0$, $e^u=\sum_{k=0}^N\frac{u^k}{k!}+r_N(u)$ with $0\le r_N(u)\le\frac{u^{N+1}}{(N+1)!}e^{u}\le\frac{e^{1/e}}{(N+1)!\,e^{N+1}}=:\rho_N$, a bound independent of $x$. Integrating over $[0,1]$,
\[
\int_0^1x^{-x}\dd x=\sum_{k=0}^N\frac1{k!}\int_0^1(-x\ln x)^k\dd x+\theta_N\rho_N,\qquad0\le\theta_N\le1.
\]
By Exercise~\ref{ex:xnex},
\[
\int_0^1(-x\ln x)^k\dd x=(-1)^k\int_0^1x^k(\ln x)^k\dd x=\frac{k!}{(k+1)^{k+1}}.
\]
So the sum is $\sum_{k=0}^N\frac1{(k+1)^{k+1}}$, that is $\sum_{n=1}^{N+1}\frac1{n^n}$, and letting $N\to\infty$ (with $\rho_N\to0$) gives the first formula. For the second, replace $u$ by $-u=x\ln x$: the same bound holds, and $\int_0^1(x\ln x)^k\dd x=\frac{(-1)^kk!}{(k+1)^{k+1}}$ gives $\sum_{n\ge1}\frac{(-1)^{n+1}}{n^n}=0.78343\ldots$ The first series converges so fast ($\frac1{10^{10}}$ at $n=10$) that ten terms give $1.2912859971$. Bernoulli published both in 1697; the nickname ``sophomore's dream'' suggests the dream that integrals and sums may be swapped at will, which here, for once, comes true.
\end{solution}

\subsection*{Notes on the challenges}

(1) Integrate the finite geometric identity. Its remainder is $R_N=(xy)^N/(1-xy)$, and on $[0,1]^2$ away from the corner,
$1-xy\ge\sqrt{(1-x)(1-y)}$. Indeed, $1-xy$ is at least each of $1-x$ and $1-y$.
Consequently
\[
0\le\int_0^1\!\!\int_0^1 R_N\dd x\dd y
\le\left(\int_0^1\frac{x^N}{\sqrt{1-x}}\dd x\right)^2
=4W_{2N+1}^2\longrightarrow0.
\]
The last equality uses $x=\sin^2t$. This bound justifies the limiting interchange independently of the series being evaluated, and gives
$\int_0^1\!\!\int_0^1(1-xy)^{-1}\dd x\dd y=\sum_{n\ge0}(n+1)^{-2}$.
After rotation the region is $0\le u\le\sqrt2$, $|v|\le\min(u,\sqrt2-u)$. Integrating in $v$ gives
$\frac4{\sqrt{2-u^2}}\arctan\frac{\min(u,\sqrt2-u)}{\sqrt{2-u^2}}$.
Put $u=\sqrt2\sin\theta$. The arctangent is $\theta$ up to $\theta=\pi/6$, and $\pi/4-\theta/2$ thereafter. Thus the integral is
\[
4\int_0^{\pi/6}\theta\dd\theta+
4\int_{\pi/6}^{\pi/2}\left(\frac\pi4-\frac\theta2\right)\dd\theta=\frac{\pi^2}6.
\]
Apostol's note is in \emph{Math.\ Intelligencer} 5 (1983), 59--60. Volume~7 supplies Fubini's theorem and the change-of-variables theorem used here.

(2) Hermite's proof that $e$ is transcendental (1873) begins with a hypothetical integer relation $\sum_{k=0}^n a_ke^k=0$, with $a_0\ne0$, and uses
$f(x)=x^{p-1}(x-1)^p\cdots(x-n)^p/(p-1)!$ for a large prime $p$.
Put $F=f+f'+\cdots+f^{(\deg f)}$. Since $(e^{-x}F(x))'=-e^{-x}f(x)$,
\[
I_k=\int_0^k e^{k-x}f(x)\dd x=e^kF(0)-F(k).
\]
The individual $I_k$ need not be integers. The relation makes their weighted sum an integer:
$\sum_{k=0}^n a_kI_k=-\sum_{k=0}^n a_kF(k)$.
For a prime $p>n$ not dividing $a_0$, the derivative values are integers, $F(k)\equiv0\pmod p$ for $1\le k\le n$, and
$F(0)\equiv(-1)^{np}(n!)^p\not\equiv0\pmod p$.
Thus this integer is nonzero. On the fixed interval $[0,n]$, the numerator of $f$ is bounded by $C^p$ for a constant $C$ independent of $p$; division by $(p-1)!$ makes the weighted integral tend to zero, a contradiction.
Lindemann's extension to $\pi$ uses $e^{i\pi}=-1$ and algebraic numbers in place of $0,\dots,n$. Hardy and Wright, Chapter~XI, gives both proofs. The transcendence of $\pi$ settles squaring the circle, since every constructible number is algebraic. Volume~3 proves the sharper fact that the degree of a constructible number is a power of $2$.

(3) If $e+\pi=r$ and $e\pi=s$ were rational, $e$ and $\pi$ would be the roots of $X^2-rX+s$, hence algebraic, contradicting Hermite and Lindemann. The irrationality of each of $e+\pi$, $e\pi$, $\pi^e$, $\pi^{\sqrt2}$ remains open; by contrast $e^\pi$ is transcendental (Gelfond 1929), and so is $e^{\pi\sqrt{163}}=262537412640768743.99999999999925\ldots$, whose nearness to an integer Volume~13 explains, from the class number of $\Q(\sqrt{-163})$.

\begin{solution}{pb:lightsout}
\idea Pressing commutes with rotating the circle, so the pressing map is a circulant matrix over $\mathbb F_2$, that is,
multiplication by a polynomial in the ring $\mathbb F_2[x]/(x^n-1)$. It is invertible exactly when that polynomial is prime
to $x^n-1$. For (b), a symmetric matrix has image orthogonal to its kernel, and the kernel is easy to guess.

\writeup Encode a set $u$ of lamps as $u(x)=\sum_{i\in u}x^i$ in $R=\mathbb F_2[x]/(x^n-1)$. Pressing the lamps of $u$
switches each lamp as many times as it has pressed lamps among itself and its neighbours, so it adds
$(x^{-1}+1+x)u=x^{-1}(1+x+x^2)u$ to the configuration (order does not matter, and pressing a lamp twice does nothing). As
$x$ is a unit in $R$, we must understand multiplication by $f=1+x+x^2$.

\emph{Lemma.} Over a field, multiplication by $f$ on $K[x]/(x^n-1)$ has a kernel of dimension $\deg g$, with
$g=\gcd(f,x^n-1)$. Indeed, $uf\equiv0$ means $x^n-1\mid uf$, that is $h\mid u\frac fg$ with $h=\frac{x^n-1}g$, and since
$\gcd(h,\frac fg)=1$ this means $h\mid u$. The multiples of $h$ of degree $<n$ are $hv$ with $\deg v<\deg g$.

(a) $1+x+x^2$ has no root in $\mathbb F_2$, so it is irreducible, and $g$ is $1$ or $1+x+x^2$. Since
$x^3-1=(x-1)(x^2+x+1)$, $x^3\equiv1$ modulo $f$, while $x\not\equiv1$ and $x^2\not\equiv1$. So $f\mid x^n-1$ if and only if
$3\mid n$. When $3\nmid n$, the kernel is $0$, pressing is a bijection of $R$, and every configuration can be switched off.
When $3\mid n$, the kernel has dimension $2$, so some configurations cannot.

(b) Let $3\mid n$, and let $v_r$ be the set of lamps $i\not\equiv r\pmod3$. Pressing $v_r$ switches each lamp twice or not at
all: a lamp in $v_r$ has exactly one neighbour in $v_r$, and a lamp outside it has both neighbours in it. So $v_0,v_1,v_2$
are in the kernel, which they span (two of them are independent and the kernel has dimension $2$). The pressing matrix
$A$ is symmetric, and over any field the image of $A$ is the orthogonal of the kernel of $A^T$ for the dot product (a
dimension count from Chapter~\ref{ch:linear}, which holds over $\mathbb F_2$ as well). So a configuration $c$ can be switched off if
and only if $c\cdot v_r=0$ for $r=0,1,2$, that is $S_1+S_2$, $S_0+S_2$, $S_0+S_1$ are all even, that is $S_0\equiv S_1\equiv S_2
\pmod2$. The image has dimension $n-2$, so this is $\frac14$ of the $2^n$ configurations.
\tools Volume~3, Chapters~6 and~9 (polynomial rings, finite fields); Chapter~\ref{ch:linear} (rank and orthogonality).
\leads The commercial game \emph{Lights Out} (1995) is the $5\times5$ grid version. Its pressing matrix has a kernel of
dimension $2$ too, so only a quarter of the boards are solvable. Sutner (1989) showed that on every finite graph pressing
all vertices is enough to reach the all-off state from all-on: a theorem about $\mathbb F_2$-linear algebra.
\end{solution}

\begin{solution}{pb:hyperplanes}
\idea If fewer than $n(q-1)$ hyperplanes did the job, the product of their equations would agree, on all of $\mathbb F_q^n$, with
a polynomial whose leading monomial is $x_1^{q-1}\cdots x_n^{q-1}$, and the Nullstellensatz forbids this.

\writeup (a) The $n(q-1)$ hyperplanes $x_j=s$, $1\le j\le n$, $s\in\mathbb F_q^*$, cover every $x\ne0$ (some $x_j\ne0$) and miss $0$.
Suppose $m<n(q-1)$ hyperplanes $H_i=\{x:a_i\cdot x=b_i\}$ cover $\mathbb F_q^n\setminus\{0\}$ and miss $0$; then $b_i\ne0$ and we
may take $b_i=1$. Let
\[
f(x)=\prod_{i=1}^m(1-a_i\cdot x)-c\prod_{j=1}^n\prod_{s\in\mathbb F_q^*}(x_j-s),
\]
with $c$ chosen so that $f(0)=0$ (the second product is non-zero at $0$). For $x\ne0$, $x$ lies on some $H_i$, so the first product
vanishes, and some $x_j\in\mathbb F_q^*$, so the second does too. Thus $f$ vanishes on $\mathbb F_q^n$. But $\deg f=n(q-1)$, since
$m<n(q-1)$, and the coefficient of $x_1^{q-1}\cdots x_n^{q-1}$ is $-c\ne0$. Problem~\ref{pb:nullstellensatz}(b) with $S_j=\mathbb F_q$ and
$t_j=q-1$ gives a point where $f\ne0$: a contradiction.

(b) Let $B$ be a blocking set; translating, we may assume $0\in B$. Every hyperplane $\{x:a\cdot x=1\}$, $a\ne0$, misses $0$, so it
contains some $y\in B\setminus\{0\}$. In other words the sets $\{a:a\cdot y=1\}$, for $y\in B\setminus\{0\}$, are hyperplanes (in the
space of the $a$'s) missing $0$ and covering every $a\ne0$. By (a), $|B|-1\ge n(q-1)$. The union of the $n$ coordinate axes, with
$1+n(q-1)$ points, is blocking: a hyperplane $a\cdot x=b$ with $b\ne0$ meets the axis $j$ for any $a_j\ne0$, and one with $b=0$
contains $0$.

\tools Problem~\ref{pb:nullstellensatz}; Chapter~\ref{ch:linear} (hyperplanes).
\leads Over $\R$ the cube $\{0,1\}^n$ minus a point needs $n$ hyperplanes, and the same proof works. Projective blocking sets are
different: in a projective plane of order $q$ a line has $q+1$ points, and the real question is how small a blocking set
containing no line can be (Blokhuis 1994: at least $\frac32(q+1)$ for $q$ prime).
\end{solution}

\begin{solution}{pb:kakeya}
\idea Linear algebra gives a low-degree polynomial vanishing on a small $K$. On each line of $K$ it vanishes at $q$ points, more
than its degree, so it vanishes identically there, and its top coefficient along the line, which is the top homogeneous part
evaluated at the direction, is $0$. A polynomial of low degree that vanishes at every direction is $0$.

\writeup (a) The polynomials of degree at most $q-1$ in $n$ variables form a space of dimension $\binom{q-1+n}n$, the number of
monomials. If $|K|<\binom{q+n-1}n$, the linear map $f\mapsto(f(x))_{x\in K}$ has a non-trivial kernel: some $f\ne0$ of degree
$d\le q-1$ vanishes on $K$. Since $K\ne\emptyset$, $f$ is not a non-zero constant, so $d\ge1$. Let $f_d\ne0$ be its homogeneous part of
degree $d$. For $v\ne0$, take $a$ with $a+tv\in K$ for all $t$. Then $t\mapsto f(a+tv)$ is a polynomial in $t$ of degree at most $d<q$
with $q$ roots, so it is $0$, and its coefficient of $t^d$, which is $f_d(v)$, is $0$. Also $f_d(0)=0$. So $f_d$ vanishes on
$\mathbb F_q^n$ while $\deg_{x_i}f_d\le q-1<q$; by Problem~\ref{pb:nullstellensatz}(a), $f_d=0$. Contradiction. Finally
$\binom{q+n-1}n=\frac{q(q+1)\cdots(q+n-1)}{n!}\ge\frac{q^n}{n!}$.

(b) For $m\in\mathbb F_q$ let $L_m$ be the line $y=mx+m^2$, of slope $m$, and let $K$ be the union of the $L_m$ and the vertical line
$x=0$. A point $(x,y)$ lies on some $L_m$ when $m^2+mx-y=0$ has a root, that is when $x^2+4y$ is a square (zero included). For fixed
$x$, $y\mapsto x^2+4y$ is a bijection of $\mathbb F_q$, and $\frac{q+1}2$ elements are squares; so $\bigcup L_m$ has $\frac{q(q+1)}2$
points. On $x=0$ these are the points $(0,y)$ with $y$ a square, $\frac{q+1}2$ of them, so the vertical line adds $\frac{q-1}2$.

\tools Chapter~\ref{ch:linear} (dimension counting); Problem~\ref{pb:nullstellensatz}(a). The construction of (b) is given here.
\leads Dvir's proof (2008) took one page and settled a problem that Wolff had posed in 1999 as a model for the real Kakeya
conjecture. Wang and Zahl proved the three-dimensional case in 2025 (arXiv:2502.17655). Dimensions $4$ and higher remain open. For $n=2$ and $q$ odd the
construction (b) is optimal (Blokhuis and Mazzocca, 2008); for $q=3$ numerical evaluation confirms by brute force that the minimum is $7$,
one more than Dvir's bound $6$.
\end{solution}

\section*{Sources and further reading}

\begin{itemize}\raggedright
\item M.~Spivak, \emph{Calculus}, 4th ed., Publish or Perish, 2008, Chapters~13--20. Integrals, Taylor's theorem and the irrationality of $e$ and $\pi$, with superb exercises.
\item E.~Hairer and G.~Wanner, \emph{Analysis by Its History}, Springer, 1996, Chapter~II. Wallis, Taylor, Lagrange's remainder, the Bernoulli--l'H\^opital contract and Euler's Gamma function, with the original texts.
\item J.~Wallis, \emph{Arithmetica Infinitorum}, Oxford, 1656; English translation with commentary by J.~A.~Stedall, \emph{The Arithmetic of Infinitesimals}, Springer, 2004. The original infinite product, found by interpolation, and the quarrel with Fermat in Stedall's introduction.
\item B.~Taylor, \emph{Methodus Incrementorum Directa et Inversa}, London, 1715, Proposition~VII; J.-L.~Lagrange, \emph{Th\'eorie des fonctions analytiques}, Paris, 1797, where the remainder appears (Part~I, \S~34--36 in the 1813 edition).
\item A.~de Moivre, \emph{Approximatio ad summam terminorum binomii $(a+b)^n$ in seriem expansi}, privately printed, 12 November 1733; reprinted in the 1738 edition of \emph{The Doctrine of Chances}. The first appearance of $e^{-x^2}$. P.-S.~Laplace, \emph{M\'emoire sur la probabilit\'e des causes par les \'ev\'enements}, 1774, for the value of the Gaussian integral.
\item C.~Truesdell, \emph{The New Bernoulli Edition}, Isis 49 (1958), 54--62, and the introduction to \emph{Der Briefwechsel von Johann Bernoulli}, vol.~1, Birkh\"auser, 1955. The l'H\^opital--Bernoulli contract of 1694 and its discovery in 1921.
\item P.~J.~Davis, \emph{Leonhard Euler's integral: a historical profile of the Gamma function}, Amer.\ Math.\ Monthly 66 (1959), 849--869. Euler's letter to Goldbach of 13 October 1729, Bohr--Mollerup, and everything else about $\Gamma$.
\item I.~Niven, \emph{A simple proof that $\pi$ is irrational}, Bull.\ Amer.\ Math.\ Soc.\ 53 (1947), 509. Exercise~\ref{ex:niven}, as it was printed. J.-B.~de Stainville, \emph{M\'elanges d'analyse alg\'ebrique et de g\'eom\'etrie}, Paris, 1815, pp.~339--341, reports Fourier's proof that $e$ is irrational (Exercise~\ref{ex:eirr}).
\item G.~H.~Hardy and E.~M.~Wright, \emph{An Introduction to the Theory of Numbers}, 6th ed., Oxford University Press, 2008, Chapter~XI. The irrationality and transcendence of $e$ and $\pi$ (Hermite, Lindemann).
\item J.~Borwein, D.~Bailey and R.~Girgensohn, \emph{Experimentation in Mathematics}, A~K Peters, 2004, Section~1.5. The sophomore's dream and Bernoulli's 1697 paper (\emph{Acta Eruditorum}, May 1697).
\item K.~Conrad, \emph{Stirling's formula} and \emph{The Gaussian integral}, expository notes on the author's web page (\url{https://kconrad.math.uconn.edu/blurbs/}). Several proofs of each, to compare with Theorem~\ref{thm:stirlingC} and Exercise~\ref{ex:gaussian}.
\item P.~Lax and M.~Terrell, \emph{Calculus with Applications}, 2nd ed., Springer, 2014. A modern calculus text with Wallis, Stirling and many numerical illustrations, including the trapezoid and Simpson rules of Proposition~\ref{prop:quadrature}.
\end{itemize}
